\documentclass[12pt,a4paper]{article}

% --- Pakiety podstawowe ---
\usepackage[polish]{babel}
\usepackage[T1]{fontenc}
\usepackage[utf8]{inputenc}
\usepackage{amsmath}
\usepackage{amsfonts}
\usepackage{amssymb}
\usepackage{graphicx}
\usepackage{geometry}
\usepackage{xcolor}
\usepackage{titlesec}
\usepackage{fancyhdr}
\usepackage{longtable}
\usepackage{booktabs}
\usepackage{array}
\usepackage[hidelinks]{hyperref}
\usepackage{enumitem}
\usepackage{calc}
\usepackage{listings}
% Pandoc-generated macro
\providecommand{\tightlist}{\setlength{\itemsep}{0pt}\setlength{\parskip}{0pt}}

% --- Konfiguracja strony ---
\geometry{margin=2.5cm}
\setlength{\parindent}{0pt}
\setlength{\parskip}{1em}

% --- Kolorystyka ---
\definecolor{NavyBlue}{RGB}{31, 56, 100}
\definecolor{MidBlue}{RGB}{46, 117, 182}
\definecolor{LightBlue}{RGB}{214, 228, 240}
\definecolor{Gray}{RGB}{242, 242, 242}

% --- Stylizacja nagłówków ---
\titleformat{\section}
  {\color{NavyBlue}\normalfont\Large\bfseries}{\thesection}{1em}{}[{\titlerule[1.5pt]}]

\titleformat{\subsection}
  {\color{MidBlue}\normalfont\large\bfseries}{\thesubsection}{1em}{}

% --- Nagłówek i stopka ---
\pagestyle{fancy}
\fancyhf{}
\fancyhead[L]{\small \color{gray} Pakiet Dokumentacji Cyberbezpieczeństwa v1.0}
\fancyhead[R]{\small \color{gray} [NAZWA FIRMY]}
\fancyfoot[C]{\thepage}
\fancyfoot[R]{\small \color{gray} POUFNE}

% --- Ustawienia tabel ---
\setlength{\LTpre}{4pt}
\setlength{\LTpost}{4pt}
\newcolumntype{L}[1]{>{\raggedright\arraybackslash}p{#1}}
\renewcommand{\arraystretch}{1.5}

\begin{document}

% --- STRONA TYTUŁOWA ---
\begin{titlepage}
    \centering
    \vspace*{2cm}
    {\huge\bfseries\color{NavyBlue} [NAZWA FIRMY KSIĘGOWEJ] \par}
    \vspace{1cm}
    {\Large\bfseries\color{MidBlue} PAKIET POLITYK I PROCEDUR\\CYBERBEZPIECZEŃSTWA \par}
    \vspace{2cm}
    
    \begin{table}[h]
        \centering
        \begin{tabular}{|>{\color{NavyBlue}\bfseries}l|l|}
            \hline
            Właściciel dokumentu: & Administrator / Właściciel firmy \\ \hline
            Status: & SZABLON -- do wypełnienia \\ \hline
            Data utworzenia: & [DATA] \\ \hline
            Następny przegląd: & [DATA + 12 MIESIĘCY] \\ \hline
            Wersja: & 1.0 \\ \hline
            Klasyfikacja: & POUFNE -- wewnętrzne \\ \hline
        \end{tabular}
    \end{table}

    \vfill
    \begin{quote}
        \centering
        \textit{\color{red} Dokument zawiera informacje poufne. Dystrybucja wyłącznie dla upoważnionych pracowników.}
    \end{quote}
    \vspace*{1cm}
\end{titlepage}

\newpage



\tableofcontents
\newpage
\phantomsection\label{doc1}
POL-001

\section{POLITYKA PRZECHOWYWANIA I ZARZĄDZANIA DANYMI
(RODO) POL-001}\label{polityka-przechowywania-i-zarzux105dzania-danymi-rodo}

\begin{longtable}[]{@{}>{\raggedright\arraybackslash}p{\dimexpr\linewidth/1000*300-2\tabcolsep\relax}>{\raggedright\arraybackslash}p{\dimexpr\linewidth/1000*700-2\tabcolsep\relax}@{}}
\toprule\noalign{}
\endhead
\bottomrule\noalign{}
\endlastfoot
Wersja: & 1.0 \\
Data wejścia w życie: & {[}DATA{]} \\
Właściciel: & Właściciel firmy / Inspektor Ochrony Danych (IOD) \\
Dotyczy: & Wszyscy pracownicy, współpracownicy, podwykonawcy \\
Przegląd: & Co 12 miesięcy lub po istotnej zmianie prawnej \\
\end{longtable}

\subsection{Cel i zakres}\label{cel-i-zakres}

Niniejsza polityka określa zasady gromadzenia, przetwarzania,
przechowywania i usuwania danych osobowych oraz danych finansowych
klientów w {[}NAZWA FIRMY{]}, zgodnie z Rozporządzeniem RODO (UE)
2016/679 oraz polskim prawem krajowym.

Polityka obejmuje: dane osobowe klientów i pracowników, dane finansowe i
księgowe, dokumenty elektroniczne i papierowe, systemy informatyczne i
nośniki danych.

\subsection{Podstawy prawne przetwarzania
danych}\label{podstawy-prawne-przetwarzania-danych}

\begin{longtable}[]{@{}>{\raggedright\arraybackslash}p{\dimexpr\linewidth/1000*270-2\tabcolsep\relax}>{\raggedright\arraybackslash}p{\dimexpr\linewidth/1000*430-2\tabcolsep\relax}>{\raggedright\arraybackslash}p{\dimexpr\linewidth/1000*300-2\tabcolsep\relax}@{}}
\toprule\noalign{}
Kategoria danych & Cel przetwarzania & Podstawa prawna (RODO) \\
\midrule\noalign{}
\endhead
\bottomrule\noalign{}
\endlastfoot
Dane klientów (osoby fizyczne) & Wykonanie usług księgowych & Art. 6
ust. 1 lit. b) i c) \\
Dane pracowników & Realizacja umowy o pracę & Art. 6 ust. 1 lit. b) i
c) \\
Dane do celów podatkowych & Obowiązek prawny (Ordynacja podatkowa) &
Art. 6 ust. 1 lit. c) \\
Dane marketingowe & Zgoda osoby & Art. 6 ust. 1 lit. a) \\
\end{longtable}

\subsection{Okresy przechowywania
danych}\label{okresy-przechowywania-danych}

\begin{longtable}[]{@{}>{\raggedright\arraybackslash}p{\dimexpr\linewidth/1000*270-2\tabcolsep\relax}>{\raggedright\arraybackslash}p{\dimexpr\linewidth/1000*430-2\tabcolsep\relax}>{\raggedright\arraybackslash}p{\dimexpr\linewidth/1000*300-2\tabcolsep\relax}@{}}
\toprule\noalign{}
Rodzaj dokumentu / danych & Okres przechowywania & Podstawa prawna \\
\midrule\noalign{}
\endhead
\bottomrule\noalign{}
\endlastfoot
Księgi rachunkowe i dokumenty finansowe & 5 lat od zakończenia roku
obrotowego & Ustawa o rachunkowości art. 74 \\
Deklaracje podatkowe & 5 lat od upływu terminu płatności podatku &
Ordynacja podatkowa art. 70 \\
Umowy z klientami & 10 lat od wygaśnięcia umowy & Kodeks cywilny --
przedawnienie \\
Akta pracownicze & 10 lat (pracownicy od 01.01.2019) & Ustawa
archiwalna \\
Dane marketingowe (zgoda) & Do cofnięcia zgody + 1 rok & RODO art. 7 \\
Logi systemowe i dostępu & 12 miesięcy & Polityka wewnętrzna \\
Kopie zapasowe (backup) & 90 dni (rotacyjnie) & Polityka wewnętrzna \\
\end{longtable}

\subsection{Prawa osób, których dane
dotyczą}\label{prawa-osuxf3b-ktuxf3rych-dane-dotyczux105}

\begin{itemize}
\tightlist
\item
  \textbf{Prawo dostępu} (art. 15 RODO) -- odpowiedź w ciągu 30 dni
\item
  \textbf{Prawo do sprostowania} (art. 16 RODO)
\item
  \textbf{Prawo do usunięcia} „prawo do bycia zapomnianym" (art. 17
  RODO)
\item
  \textbf{Prawo do ograniczenia przetwarzania} (art. 18 RODO)
\item
  \textbf{Prawo do przenoszenia danych} (art. 20 RODO)
\item
  \textbf{Prawo do sprzeciwu} (art. 21 RODO)
\end{itemize}

Punkt kontaktowy dla wniosków: \textbf{{[}EMAIL/ADRES{]}} \textbar{}
Termin rozpatrzenia: 30 dni (możliwość przedłużenia o 60 dni).

\subsection{Bezpieczeństwo danych}\label{bezpieczeux144stwo-danych}

\begin{itemize}
\tightlist
\item
  Dane przetwarzane wyłącznie na zatwierdzonych urządzeniach i
  platformach
\item
  Szyfrowanie danych w spoczynku (AES-256) i podczas transmisji (TLS
  1.2+)
\item
  Dostęp na zasadzie „najmniejszych uprawnień" (least privilege)
\item
  Regularne kopie zapasowe wg polityki backup
\item
  Pseudonimizacja danych tam, gdzie to możliwe
\end{itemize}

\subsection{Naruszenie ochrony danych --
procedura}\label{naruszenie-ochrony-danych-procedura}

\begin{enumerate}
\tightlist
\item
  Wykrycie naruszenia → natychmiastowe poinformowanie przełożonego /
  właściciela
\item
  Ocena ryzyka -- czy naruszenie stwarza ryzyko dla praw osób?
\item
  \textbf{Zgłoszenie do UODO} -- w ciągu \textbf{72 godzin} od wykrycia
  (jeśli ryzyko wystąpiło)
\item
  Powiadomienie osób poszkodowanych -- bez zbędnej zwłoki (jeśli wysokie
  ryzyko)
\item
  Dokumentacja incydentu -- rejestr naruszeń (wewnętrzny)
\end{enumerate}

\subsection{Powierzenie danych podmiotom
zewnętrznym}\label{powierzenie-danych-podmiotom-zewnux119trznym}

Każde powierzenie danych wymaga zawarcia \textbf{umowy powierzenia
przetwarzania (DPA)}. Lista zatwierdzonych procesorów:

\begin{longtable}[]{@{}>{\raggedright\arraybackslash}p{\dimexpr\linewidth/1000*270-2\tabcolsep\relax}>{\raggedright\arraybackslash}p{\dimexpr\linewidth/1000*430-2\tabcolsep\relax}>{\raggedright\arraybackslash}p{\dimexpr\linewidth/1000*300-2\tabcolsep\relax}@{}}
\toprule\noalign{}
Procesor & Zakres przetwarzania & Umowa DPA \\
\midrule\noalign{}
\endhead
\bottomrule\noalign{}
\endlastfoot
{[}NAZWA DOSTAWCY CHMURY{]} & Przechowywanie plików & {[} {]} TAK ~{[}
{]} NIE ~Data: \underline{\hspace{2cm}} \\
{[}BIURO IT / ZEWNĘTRZNY IT{]} & Wsparcie informatyczne & {[} {]} TAK
~{[} {]} NIE ~Data: \underline{\hspace{2cm}} \\
{[}INNE{]} & \underline{\hspace{3cm}} & {[} {]} TAK ~{[}
{]} NIE ~Data: \underline{\hspace{2cm}} \\
\end{longtable}

\begin{longtable}[]{@{}
  >{\raggedright\arraybackslash}p{\dimexpr\linewidth/1000*500-1\tabcolsep\relax}
  >{\raggedright\arraybackslash}p{\dimexpr\linewidth/1000*500-1\tabcolsep\relax}@{}}
\toprule\noalign{}
\endhead
\bottomrule\noalign{}
\endlastfoot
\begin{minipage}[t]{\linewidth}\raggedright
Zatwierdził:

Podpis: \underline{\hspace{3cm}}

Data: \underline{\hspace{4cm}}
\end{minipage} & \begin{minipage}[t]{\linewidth}\raggedright
Stanowisko / Tytuł:

Imię i nazwisko: \underline{\hspace{3cm}}

Funkcja: \underline{\hspace{3cm}}
\end{minipage} \\
\end{longtable}

\phantomsection\label{doc2}
POL-002

\section{POLITYKA DOSTĘPU I MONITOROWANIA
UŻYTKOWNIKÓW}\label{polityka}

\begin{longtable}[]{@{}>{\raggedright\arraybackslash}p{\dimexpr\linewidth/1000*300-2\tabcolsep\relax}>{\raggedright\arraybackslash}p{\dimexpr\linewidth/1000*700-2\tabcolsep\relax}@{}}
\toprule\noalign{}
\endhead
\bottomrule\noalign{}
\endlastfoot
Wersja: & 1.0 \\
Dotyczy: & Wszyscy użytkownicy systemów informatycznych firmy \\
Odpowiedzialny: & Administrator IT / Kierownik \\
Przegląd: & Co 12 miesięcy lub przy zmianie struktury organizacyjnej \\
Norma odniesienia: & ISO/IEC 27001:2022, Załącznik A - kontrole: A.5.15, A.5.16, A.5.17, A.5.18, A.8.2, A.8.3, A.8.5, A.5.10, A.8.15 \\
\end{longtable}

\subsection{Cel i zakres polityki}\label{cel-i-zakres-pol}

Niniejsza polityka określa zasady zarządzania dostępem użytkowników do systemów informatycznych firmy [NAZWA FIRMY KSIĘGOWEJ], monitorowania działań użytkowników oraz dopuszczalnego korzystania z zasobów IT. Polityka realizuje wymagania normy ISO/IEC 27001:2022 w zakresie kontroli dostępu (kontrole A.5.15-A.5.18, A.8.2-A.8.5)

Zakres obejmuje: konta użytkowników w systemach operacyjnych, system księgowy, pocztę e-mail, VPN, zasoby sieciowe (serwer plików), usługi chmurowe, urządzenia końcowe (laptopy, stacje robocze).

\subsection{Polityka kontroli dostępu}\label{polityka-kontroli-dostepu}

Zgodnie z kontrolą A.5.15 (Kontrola dostępu):

\begin{itemize}
\tightlist
\item
  Dostęp do systemów informatycznych jest przyznawany wyłącznie osobom posiadającym uzasadnioną potrzebę biznesową
\item
  Dostęp wymaga formalnej autoryzacji - pisemny wniosek kierownika lub właściciela firmy
\item
  Każdy użytkownik posiada \textbf{indywidualne, unikalne konto} -współdzielenie kont jest zabronione
\item
  Konta usługowe (\textit{service accounts}) posiadają wyznaczonego właściciela i podlegają tym samym zasadom przeglądu co konta użytkowników
\end{itemize}

\subsection{Zasada najniższych uprawnień i rozdzielenia obowiązków}\label{zasada}

Zgodnie z kontrolami A.5.15 (Kontrola dostępu) i A.8.2 (Uprzywilejowane prawa dostępu)

\begin{itemize}
\tightlist
\item
  Każdy użytkownik otrzymuje dostęp wyłącznie do zasobów \textbf{niezbędnych} do wykonywania swoich obowiązków (\textit{least privilege})
\item
  Uprawnienia administracyjne (konta uprzywilejowane) przydzielane wyłącznie administratorowi IT i właścicielowi firmy
\item
  Konta uprzywilejowane \textbf{nie są używane} do codziennej pracy - administrator posiada osobne konto standardowe
\item
  Rozdzielenie obowiązków (\textit{segregation of duties}): jedna osoba nie może jednocześnie wprowadzać i zatwierdzać operacji finansowych w systemie księgowym
\item
  Uprawnienia są nadawane przez administratora IT na podstawie pisemnego wniosku zaakceptowanego przez kierownika/właściciela
\end{itemize}
\newpage
\subsection{Macierz dostępu}\label{macierz-dostepu-szablon}

Zgodnie z kontrolą A.8.3 (Ograniczenie dostępu do informacji):

\begin{longtable}[]{@{}>{\raggedright\arraybackslash}p{\dimexpr\linewidth/1000*220-2\tabcolsep\relax}>{\raggedright\arraybackslash}p{\dimexpr\linewidth/1000*200-2\tabcolsep\relax}>{\raggedright\arraybackslash}p{\dimexpr\linewidth/1000*200-2\tabcolsep\relax}>{\raggedright\arraybackslash}p{\dimexpr\linewidth/1000*200-2\tabcolsep\relax}>{\raggedright\arraybackslash}p{\dimexpr\linewidth/1000*180-2\tabcolsep\relax}@{}}
\toprule\noalign{}
Zasób / System & Właściciel & Księgowy & Asystent & Gość/Klient \\
\midrule\noalign{}
\endhead
\bottomrule\noalign{}
\endlastfoot
System księgowy (pełny) & Pełny & Pełny & Odczyt & Brak \\
Dane klientów (CRM) & Pełny & Odczyt/Zapis & Odczyt & Brak \\
Dokumenty finansowe (serwer) & Pełny & Swoje teczki & Brak & Brak \\
Poczta e-mail firmowa & Pełny & Własna skrzynka & Własna skrzynka & Brak \\
Internet / VPN & Pełny & Tak & Ograniczony & Brak \\
Panel administracyjny IT & Pełny & Brak & Brak & Brak \\
Kopie zapasowe & Pełny & Brak & Brak & Brak \\
Logi i narzędzia audytowe & Pełny & Brak & Brak & Brak \\
\end{longtable}

Macierz dostępu podlega przeglądowi co \textbf{kwartał} (kontrola A.5.18)

\subsection{Zarządzanie cyklem życia konta użytkownika}\label{cykl-zycia-konta}

Zgodnie z kontrolami A.5.16 (Zarządzanie tożsamością) i A.5.18 (Prawa dostępu).

\begin{longtable}[]{@{}>{\raggedright\arraybackslash}p{\dimexpr\linewidth/1000*250-2\tabcolsep\relax}>{\raggedright\arraybackslash}p{\dimexpr\linewidth/1000*400-2\tabcolsep\relax}>{\raggedright\arraybackslash}p{\dimexpr\linewidth/1000*350-2\tabcolsep\relax}@{}}
\toprule\noalign{}
Etap & Działanie & Odpowiedzialny \\
\midrule\noalign{}
\endhead
\bottomrule\noalign{}
\endlastfoot
Utworzenie konta & Wniosek pisemny + nadanie uprawnień wg macierzy & Admin IT (na wniosek kierownika) \\
Zmiana roli/stanowiska & Przegląd i aktualizacja uprawnień -- usunięcie zbędnych, nadanie nowych & Admin IT (na wniosek kierownika) \\
Długa nieobecność ($>$30 dni) & Tymczasowa dezaktywacja konta & Admin IT (powiadomienie od HR) \\
Zakończenie współpracy & Natychmiastowa dezaktywacja konta -- procedura PROC-007 & Admin IT + HR \\
Archiwizacja & Po 90 dniach od dezaktywacji -- usunięcie konta, archiwizacja danych & Admin IT \\
Przegląd kwartalny & Weryfikacja aktywnych kont i uprawnień z macierzą & Admin IT \\
\end{longtable}

\subsection{Zarządzanie poświadczeniami użytkowników}\label{zarzadzanie-poswiadczeniami}

Zgodnie z kontrolą A.5.17 (Informacje uwierzytelniające):

\begin{itemize}
\tightlist
\item
  Hasła początkowe generowane przez administratora IT -użytkownik zobowiązany do zmiany przy pierwszym logowaniu
\item
  Hasło tymczasowe przekazywane wyłącznie w sposób bezpieczny (osobiście lub zaszyfrowanym kanałem) - nigdy w zwykłym e-mailu
\item
  Wymagania dotyczące haseł określone w PROC-004 (sekcja 4: Hasła i uwierzytelnianie)
\item
  Obowiązkowe MFA dla wszystkich kont z dostępem do danych klientów lub systemów księgowych
\item
  Użytkownik ponosi odpowiedzialność za zachowanie poufności swoich poświadczeń
\item
  W przypadku podejrzenia kompromitacji poświadczeń -natychmiastowe zgłoszenie do administratora IT i zmiana hasła
\end{itemize}

\subsection{Monitorowanie działań użytkowników}\label{monitorowanie}

Zgodnie z kontrolami A.8.15 (Logowanie) i A.8.16 (Działania monitorujące) oraz zgodnie z art. 22 Kodeksu pracy:

Firma zastrzega sobie prawo do monitorowania korzystania z systemów informatycznych w celu zapewnienia bezpieczeństwa informacji i ochrony zasobów firmy. Użytkownicy są poinformowani o fakcie, zakresie i celu monitorowania (klauzula informacyjna w umowie o pracę, regulaminie IT oraz przy onboardingu - PROC-006).

\subsubsection{Zakres monitorowania}\label{zakres-monitorowania}

\begin{longtable}[]{@{}>{\raggedright\arraybackslash}p{\dimexpr\linewidth/1000*350-2\tabcolsep\relax}>{\raggedright\arraybackslash}p{\dimexpr\linewidth/1000*350-2\tabcolsep\relax}>{\raggedright\arraybackslash}p{\dimexpr\linewidth/1000*300-2\tabcolsep\relax}@{}}
\toprule\noalign{}
Monitorowane działanie & Cel monitorowania & Kontrola ISO \\
\midrule\noalign{}
\endhead
\bottomrule\noalign{}
\endlastfoot
Logowania i wylogowania (sukces/porażka) & Wykrywanie nieautoryzowanego dostępu & A.8.15 \\
Dostęp do plików i baz danych & Ochrona poufności danych klientów & A.8.3, A.8.15 \\
Wysyłane e-maile z firmowych kont (metadane) & Zapobieganie wyciekowi danych & A.8.15, A.5.14 \\
Dostęp do Internetu (adresy URL) & Ochrona przed złośliwym oprogramowaniem & A.8.23 \\
Drukowanie dokumentów & Kontrola obiegu dokumentów poufnych & A.8.15 \\
Podłączanie urządzeń zewnętrznych (USB) & Ochrona przed wyciekiem danych i malware & A.8.7, A.8.12 \\
Operacje na kontach uprzywilejowanych & Rozliczalność działań administratora & A.8.2, A.8.15 \\
\end{longtable}

\subsubsection{Ochrona danych z monitorowania}\label{ochrona-danych-monitoring}

\begin{itemize}
\tightlist
\item
  Dane z monitorowania przechowywane przez \textbf{12 miesięcy} (logi systemowe) lub \textbf{5 lat} (logi systemu księgowego)
\item
  Dostęp do danych z monitorowania: wyłącznie administrator IT i właściciel firmy
\item
  Dane z monitorowania zawierające dane osobowe podlegają ochronie zgodnie z RODO - podstawą przetwarzania jest uzasadniony interes administratora (art. 6 ust. 1 lit. f RODO)
\item
  Szczegółowe zasady przechowywania i ochrony logów określa POL-009
\end{itemize}

\subsection{Akceptowalne użycie zasobów IT}\label{akceptowalne}

Zgodnie z kontrolą A.5.10 (Akceptowalne użycie informacji i innych powiązanych aktywów):

\subsubsection{Dozwolone:}\label{dozwolone}

\begin{itemize}
\tightlist
\item
  Używanie sprzętu i oprogramowania firmy do celów służbowych
\item
  Ograniczone użycie prywatne niezakłócające pracy i niepowodujące zagrożeń bezpieczeństwa (np. krótkie przeglądanie wiadomości w przerwie)
\item
  Korzystanie wyłącznie z zatwierdzonych aplikacji i usług (lista w PROC-003)
\end{itemize}

\subsubsection{ Zabronione:}\label{zabronione}

\begin{itemize}
\tightlist
\item
  Instalacja nieautoryzowanego oprogramowania (bez zgody administratora IT)
\item
  Podłączanie prywatnych nośników USB bez pisemnej zgody administratora
\item
  Udostępnianie danych logowania (hasła, tokeny MFA) innym osobom -- w tym współpracownikom
\item
  Obchodzenie zabezpieczeń: wyłączanie antywirusa/EDR, omijanie VPN, modyfikacja ustawień firewall
\item
  Przesyłanie danych klientów na prywatne skrzynki e-mail, prywatne usługi chmurowe lub niezatwierdzone nośniki
\item
  Używanie służbowych kont e-mail do celów prywatnych lub rejestracji w serwisach niezwiązanych z pracą
\item
  Kopiowanie danych firmowych na urządzenia prywatne bez autoryzacji
\item
  Uzyskiwanie dostępu do zasobów, do których użytkownik nie posiada uprawnień (nawet w celach testowych)
\end{itemize}

\subsection{Konsekwencje naruszenia polityki}\label{konsekwencje-naruszenia}

Naruszenie niniejszej polityki może skutkować:

\begin{itemize}
\tightlist
\item
  Natychmiastowym ograniczeniem lub cofnięciem uprawnień dostępu
\item
  Postępowaniem dyscyplinarnym zgodnie z regulaminem pracy
\item
  Odpowiedzialnością karną w przypadku naruszenia przepisów o ochronie danych osobowych
\item
  Odpowiedzialnością cywilną za szkody wyrządzone firmie lub klientom
\end{itemize}

Każdy incydent naruszenia polityki jest dokumentowany i analizowany przez administratora IT i właściciela firmy.

\subsection{Przegląd uprawnień - harmonogram}\label{przeglad-uprawnien-harmonogram}

Zgodnie z kontrolą A.5.18 (Prawa dostępu):

\begin{longtable}[]{@{}>{\raggedright\arraybackslash}p{\dimexpr\linewidth/1000*350-2\tabcolsep\relax}>{\raggedright\arraybackslash}p{\dimexpr\linewidth/1000*200-2\tabcolsep\relax}>{\raggedright\arraybackslash}p{\dimexpr\linewidth/1000*200-2\tabcolsep\relax}>{\raggedright\arraybackslash}p{\dimexpr\linewidth/1000*250-2\tabcolsep\relax}@{}}
\toprule\noalign{}
Rodzaj przeglądu & Częstotliwość & Odpowiedzialny & Ostatni / Następny \\
\midrule\noalign{}
\endhead
\bottomrule\noalign{}
\endlastfoot
Weryfikacja aktywnych kont & Co miesiąc & Admin IT & \underline{\hspace{2cm}} / \underline{\hspace{2cm}} \\
Przegląd uprawnień (porównanie z macierzą) & Co kwartał & Admin IT & \underline{\hspace{2cm}} / \underline{\hspace{2cm}} \\
Przegląd kont uprzywilejowanych & Co kwartał & Właściciel + Admin IT & \underline{\hspace{2cm}} / \underline{\hspace{2cm}} \\
Przegląd kont usługowych & Co 6 miesięcy & Admin IT & \underline{\hspace{2cm}} / \underline{\hspace{2cm}} \\
Pełny audyt polityki dostępu & Co 12 miesięcy & Właściciel + Admin IT & \underline{\hspace{2cm}} / \underline{\hspace{2cm}} \\
\end{longtable}

\subsection{Rejestr zmian i przeglądów polityki}\label{rejestr-zmian}

\begin{longtable}[]{@{}>{\raggedright\arraybackslash}p{\dimexpr\linewidth/1000*150-2\tabcolsep\relax}>{\raggedright\arraybackslash}p{\dimexpr\linewidth/1000*150-2\tabcolsep\relax}>{\raggedright\arraybackslash}p{\dimexpr\linewidth/1000*450-2\tabcolsep\relax}>{\raggedright\arraybackslash}p{\dimexpr\linewidth/1000*250-2\tabcolsep\relax}@{}}
\toprule\noalign{}
Wersja & Data & Opis zmiany & Autor \\
\midrule\noalign{}
\endhead
\bottomrule\noalign{}
\endlastfoot
1.0 & {[}DATA{]} & Wersja początkowa & {[}IMIĘ NAZWISKO{]} \\
\underline{\hspace{1cm}} & \underline{\hspace{2cm}} & \underline{\hspace{5cm}} & \underline{\hspace{3cm}} \\
\underline{\hspace{1cm}} & \underline{\hspace{2cm}} & \underline{\hspace{5cm}} & \underline{\hspace{3cm}} \\
\end{longtable}

\begin{longtable}[]{@{}
  >{\raggedright\arraybackslash}p{\dimexpr\linewidth/1000*500-1\tabcolsep\relax}
  >{\raggedright\arraybackslash}p{\dimexpr\linewidth/1000*500-1\tabcolsep\relax}@{}}
\toprule\noalign{}
\endhead
\bottomrule\noalign{}
\endlastfoot
\begin{minipage}[t]{\linewidth}\raggedright
Zatwierdził:

Podpis: \underline{\hspace{3cm}}

Data: \underline{\hspace{4cm}}
\end{minipage} & \begin{minipage}[t]{\linewidth}\raggedright
Stanowisko / Tytuł:

Imię i nazwisko: \underline{\hspace{3cm}}

Funkcja: \underline{\hspace{3cm}}
\end{minipage} \\
\end{longtable}
\phantomsection\label{doc6}
PROC-003

\section{PROCEDURA ONBOARDINGU NOWYCH UŻYTKOWNIKÓW
(IT) PROC-003}\label{procedura-onboardingu-nowych-uux17cytkownikuxf3w-it}

\begin{longtable}[]{@{}>{\raggedright\arraybackslash}p{\dimexpr\linewidth/1000*300-2\tabcolsep\relax}>{\raggedright\arraybackslash}p{\dimexpr\linewidth/1000*700-2\tabcolsep\relax}@{}}
\toprule\noalign{}
\endhead
\bottomrule\noalign{}
\endlastfoot
Wersja: & 1.0 \\
Dotyczy: & Nowi pracownicy, stażyści, współpracownicy \\
Odpowiedzialny: & Administrator IT + HR \\
Termin realizacji: & Dzień przed lub w dniu pierwszego dnia pracy \\
\end{longtable}
\newpage
\subsection{Lista kontrolna onboardingu
IT}\label{lista-kontrolna-onboardingu-it}

\begin{longtable}[]{@{}>{\raggedright\arraybackslash}p{\dimexpr\linewidth/1000*270-2\tabcolsep\relax}>{\raggedright\arraybackslash}p{\dimexpr\linewidth/1000*250-2\tabcolsep\relax}>{\raggedright\arraybackslash}p{\dimexpr\linewidth/1000*250-2\tabcolsep\relax}>{\raggedright\arraybackslash}p{\dimexpr\linewidth/1000*230-2\tabcolsep\relax}@{}}
\toprule\noalign{}
\# & Czynność & Odpowiedzialny & Status / Data \\
\midrule\noalign{}
\endhead
\bottomrule\noalign{}
\endlastfoot
1 & Odebranie wypełnionego wniosku o konto od kierownika & HR /
Kierownik & {[} {]} \underline{\hspace{2cm}} \\
2 & Utworzenie konta użytkownika w AD / systemie & Administrator IT &
{[} {]} \underline{\hspace{2cm}} \\
3 & Nadanie uprawnień wg macierzy dostępu (POL-005) & Administrator IT &
{[} {]} \underline{\hspace{2cm}} \\
4 & Konfiguracja MFA (telefon / aplikacja) & Administrator IT & {[} {]}
\underline{\hspace{2cm}} \\
5 & Wydanie sprzętu -- laptop, akcesoria (protokół zdawczo-odbiorczy) &
Administrator IT & {[} {]} \underline{\hspace{2cm}} \\
6 & Instalacja i konfiguracja wymaganego oprogramowania & Administrator
IT & {[} {]} \underline{\hspace{2cm}} \\
7 & Konfiguracja poczty e-mail i podpisu & Administrator IT & {[} {]}
\underline{\hspace{2cm}} \\
8 & Podłączenie do VPN + test połączenia & Administrator IT & {[} {]}
\underline{\hspace{2cm}} \\
9 & Zapoznanie z Polityką Bezpieczeństwa IT (POL-001, POL-005) &
Administrator IT & {[} {]} \underline{\hspace{2cm}} \\
10 & Szkolenie wprowadzające z cyberbezpieczeństwa (PROC-008) &
Administrator IT & {[} {]} \underline{\hspace{2cm}} \\
11 & Podpisanie Oświadczenia o zapoznaniu z politykami IT & HR /
Pracownik & {[} {]} \underline{\hspace{2cm}} \\
12 & Podpisanie Oświadczenia o poufności (NDA / klauzula RODO) & HR /
Pracownik & {[} {]} \underline{\hspace{2cm}} \\
13 & Test dostępu -- potwierdzenie że wszystko działa & Nowy pracownik &
{[} {]} \underline{\hspace{2cm}} \\
\end{longtable}

\subsection{Oświadczenie o zapoznaniu z politykami IT (do
podpisania)}\label{oux15bwiadczenie-o-zapoznaniu-z-politykami-it-do-podpisania}

Ja, niżej podpisany/a {~}, potwierdzam, że:

\begin{enumerate}
\tightlist
\item
  Zapoznałem/am się z Polityką Bezpieczeństwa IT firmy {[}NAZWA{]}
\item
  Zostałem/am poinformowany/a o fakcie monitorowania aktywności w
  systemach IT
\item
  Zobowiązuję się do przestrzegania zasad określonych w politykach
\item
  Przyjmuję do wiadomości, że naruszenie polityk może skutkować
  konsekwencjami służbowymi
\end{enumerate}

Data: \underline{\hspace{3cm}}

Podpis pracownika: \underline{\hspace{3cm}}

Podpis przełożonego: \underline{\hspace{3cm}}

\begin{longtable}[]{@{}
  >{\raggedright\arraybackslash}p{\dimexpr\linewidth/1000*500-1\tabcolsep\relax}
  >{\raggedright\arraybackslash}p{\dimexpr\linewidth/1000*500-1\tabcolsep\relax}@{}}
\toprule\noalign{}
\endhead
\bottomrule\noalign{}
\endlastfoot
\begin{minipage}[t]{\linewidth}\raggedright
Zatwierdził:

Podpis: \underline{\hspace{3cm}}

Data: \underline{\hspace{4cm}}
\end{minipage} & \begin{minipage}[t]{\linewidth}\raggedright
Stanowisko / Tytuł:

Imię i nazwisko: \underline{\hspace{3cm}}

Funkcja: \underline{\hspace{3cm}}
\end{minipage} \\
\end{longtable}

\phantomsection\label{doc7}
PROC-004

\section{PROCEDURA OFFBOARDINGU ODCHODZĄCYCH UŻYTKOWNIKÓW
(IT) PROC-004}\label{procedura-offboardingu-odchodzux105cych-uux17cytkownikuxf3w-it}

\begin{longtable}[]{@{}>{\raggedright\arraybackslash}p{\dimexpr\linewidth/1000*300-2\tabcolsep\relax}>{\raggedright\arraybackslash}p{\dimexpr\linewidth/1000*700-2\tabcolsep\relax}@{}}
\toprule\noalign{}
\endhead
\bottomrule\noalign{}
\endlastfoot
Wersja: & 1.0 \\
Dotyczy: & Pracownicy, stażyści, współpracownicy kończący współpracę \\
Odpowiedzialny: & Administrator IT + HR \\
Termin realizacji: & W dniu lub natychmiast po ostatnim dniu pracy \\
\end{longtable}

\textbf{UWAGA:} Konta należy dezaktywować w dniu zakończenia współpracy -- nie
wcześniej, nie później (ryzyko wycieku po odejściu lub brak dostępu w
ostatnim dniu).

\subsection{Lista kontrolna offboardingu
IT}\label{lista-kontrolna-offboardingu-it}

\begin{longtable}[]{@{}>{\raggedright\arraybackslash}p{\dimexpr\linewidth/1000*270-2\tabcolsep\relax}>{\raggedright\arraybackslash}p{\dimexpr\linewidth/1000*250-2\tabcolsep\relax}>{\raggedright\arraybackslash}p{\dimexpr\linewidth/1000*250-2\tabcolsep\relax}>{\raggedright\arraybackslash}p{\dimexpr\linewidth/1000*230-2\tabcolsep\relax}@{}}
\toprule\noalign{}
\# & Czynność & Odpowiedzialny & Status / Data \\
\midrule\noalign{}
\endhead
\bottomrule\noalign{}
\endlastfoot
1 & Wyłączenie konta użytkownika (nie usuwanie -- archiwum 90 dni) &
Administrator IT & {[} {]} \underline{\hspace{2cm}} \\
2 & Zmiana haseł do współdzielonych kont (jeśli dotyczy) & Administrator
IT & {[} {]} \underline{\hspace{2cm}} \\
3 & Odwołanie dostępu VPN & Administrator IT & {[} {]}
\underline{\hspace{2cm}} \\
4 & Usunięcie / dezaktywacja tokena MFA & Administrator IT & {[} {]}
\underline{\hspace{2cm}} \\
5 & Archiwizacja skrzynki e-mail (90 dni przekierowania do przełożonego)
& Administrator IT & {[} {]} \underline{\hspace{2cm}} \\
6 & Przejęcie plików służbowych (przekazanie kierownikowi) & Kierownik &
{[} {]} \underline{\hspace{2cm}} \\
7 & Zwrot sprzętu (laptop, telefon, klucze, karty dostępu) -- protokół &
HR & {[} {]} \underline{\hspace{2cm}} \\
8 & Wyczyszczenie danych z urządzeń zwracanych (certyfikowane usunięcie)
& Administrator IT & {[} {]} \underline{\hspace{2cm}} \\
9 & Usunięcie z list mailingowych i grup & Administrator IT & {[} {]}
\underline{\hspace{2cm}} \\
10 & Odwołanie dostępu do usług zewnętrznych (chmura, narzędzia SaaS) &
Administrator IT & {[} {]} \underline{\hspace{2cm}} \\
11 & Podpisanie oświadczenia o zwrocie danych i poufności & HR /
Pracownik & {[} {]} \underline{\hspace{2cm}} \\
12 & Aktualizacja macierzy dostępu i inwentarza sprzętowego &
Administrator IT & {[} {]} \underline{\hspace{2cm}} \\
\end{longtable}

\subsection{Oświadczenie o zwrocie danych i poufności
(Offboarding)}\label{oux15bwiadczenie-o-zwrocie-danych-i-poufnoux15bci-offboarding}

Ja, niżej podpisany/a {~}, potwierdzam, że:

\begin{enumerate}
\tightlist
\item
  Zwróciłem/am cały sprzęt firmowy (wymieniony w załączniku protokołu)
\item
  Usunąłem/am wszelkie dane firmowe z urządzeń prywatnych
\item
  Nie posiadam kopii danych klientów, dokumentów finansowych ani haseł
  firmowych
\item
  Zobowiązuję się do zachowania poufności przez {[}3/5{]} lat po
  zakończeniu współpracy
\item
  Jestem świadomy/a konsekwencji prawnych naruszenia powyższych
  zobowiązań
\end{enumerate}

Data: \underline{\hspace{3cm}}

Podpis pracownika: \underline{\hspace{3cm}}

Podpis HR / przełożonego: \underline{\hspace{3cm}}

\begin{longtable}[]{@{}
  >{\raggedright\arraybackslash}p{\dimexpr\linewidth/1000*500-1\tabcolsep\relax}
  >{\raggedright\arraybackslash}p{\dimexpr\linewidth/1000*500-1\tabcolsep\relax}@{}}
\toprule\noalign{}
\endhead
\bottomrule\noalign{}
\endlastfoot
\begin{minipage}[t]{\linewidth}\raggedright
Zatwierdził:

Podpis: \underline{\hspace{3cm}}

Data: \underline{\hspace{4cm}}
\end{minipage} & \begin{minipage}[t]{\linewidth}\raggedright
Stanowisko / Tytuł:

Imię i nazwisko: \underline{\hspace{3cm}}

Funkcja: \underline{\hspace{3cm}}
\end{minipage} \\
\end{longtable}

\newpage
\phantomsection\label{doc8}
PROC-005

\section{PROCEDURA WERYFIKACJI CERTYFIKACJI I BEZPIECZEŃSTWA
SPRZĘTU PROC-005}\label{procedura-weryfikacji-certyfikacji-i-bezpieczeux144stwa-sprzux119tu}

\begin{longtable}[]{@{}>{\raggedright\arraybackslash}p{\dimexpr\linewidth/1000*300-2\tabcolsep\relax}>{\raggedright\arraybackslash}p{\dimexpr\linewidth/1000*700-2\tabcolsep\relax}@{}}
\toprule\noalign{}
\endhead
\bottomrule\noalign{}
\endlastfoot
Wersja: & 1.0 \\
Dotyczy: & Komputery, laptopy, switche, routery, drukarki, UPS \\
Odpowiedzialny: & Administrator IT / Właściciel \\
Częstotliwość: & Przy zakupie + coroczna weryfikacja \\
\end{longtable}

\subsection{ Wymagania przy zakupie nowego
sprzętu}\label{wymagania-przy-zakupie-nowego-sprzux119tu}

\subsubsection{Komputery i laptopy}\label{komputery-i-laptopy}
Komputery i laptopy należą do grupy urządzeń udostępnianych wszystkim pracownikom organizacji. Sprawia to, że są jednymi z najbardziej narażonych elementów infrastruktury komputerowej, co atakujący mogą próbować wykorzystać. W celu zapewnienia jak najwyższego poziomu bezpieczeństwa określono poniższe wymagania, które powinny zostać spełnione podczas zarówno zakupu jak i eksploatacji oraz utylizacji urządzeń:
\begin{enumerate}
    \item Zakup komputerów i laptopów Sprzęt, niezależnie od swojej złożoności czy wartości nominalnej, powinien zostać zakupiony od autoryzowanego dystrybutora. Zakup powinien zostać udokumentowany w postaci faktury VAT. Sprzęt powinien być objęty gwarancją producenta oraz pochodzić od renomowanej marki. Nie dopuszcza się zakupu urządzeń używanych objętych gwarancją producenta (B-stock) ani urządzeń pochodzących z rynku wtórnego (second-hand).
    
    \item Wymagane certyfikaty Zakupione jednostki powinna posiadać odpowiednią certyfikację producenta. W celu zapewnienia odpowiedniego poziomu bezpieczeństwa sprzętu, wyszczególniono poniższe certyfikaty:
        \begin{itemize}
            \item Certyfikat Europejski (CE) - podstawowy certyfikat, który musi spełniać każdy produkt sprzedawany na terenie Unii Europejskiej.
            \item Certyfikat Federalnej Komisji Łączności (FCC) - certyfikat zapewniający, że zakłócenia elektromagnetyczne generowane przez urządzenie mieszczą się w granicach normy określonej przez Federalną Komisję Łączności
            \item Certyfikat RoHS - certyfikat zapewniający, że urządzenie nie składa się z nadmiernej ilości substancji niebezpiecznych dla zdrowia
        \end{itemize}

    \item Wymagania minimalne Jednostki komputerowe powinny spełniać minimalne wymagania konstrukcyjne oraz systemowe jakimi są:
        \begin{itemize}
            \item TPM 2.0 - układ scalony przechowujący klucze enkrypcji do aplikacji czy dysków twardych
            \item Secure Boot - oprogramowanie zapobiegające uruchamianiu się złośliwego oprogramowania przed rozruchem systemu
            \item Szyfrowanie Dysku (BitLocker/FileValult)
        \end{itemize}

    \item Zatwierdzenie przez administratora IT Zakup sprzętu komputerowego powinien każdorazowo zostać zatwierdzony przez administratora IT po wcześniejszym zweryfikowaniu potrzeb oraz dostępnych zasobów sprzętowych organizacji. Ograniczy to nadmiarową ilość sprzętu oraz pozwoli na lepsze dostosowanie jednostek do zapotrzebowania firmy.
    
    \item Rejestracja w inwentarzu sprzętowym Po zakupieniu, jednostki powinny zostać zarejestrowane w inwentarzu sprzętowym. Wpis jednostki powinien zawierać datę zakupu, jej numer seryjny, numer faktury oraz obecnego użytkownika do którego ją przypisano.
\end{enumerate}

\subsubsection{Routery i switche}\label{routery-i-switche}
Routery oraz switche są jednymi z kluczowych komponentów infrastruktury sieciowej organizacji. Podobnie jak komputery, powinny one spełniać szereg wymogów w celu zapewnienia odpowiedniego poziomu bezpieczeństwa całej infrastruktury.

\begin{enumerate}
    \item Zakup od renomowanych producentów Tak samo jak w przypadku komputerów i laptopów, zaleca się zakup sprzętu od renomowanych producentów, takich jak Cisco, Fortinet, Ubiquiti czy TP-Link. Sprzęt powinien zostać zakupiony od autoryzowanego dystrybutora jako nowy. Nie dopuszcza się zakupu sprzętu używanego objętego gwarancją producenta (B-stock) ani sprzętu z rynku wtórnego (second-hand).
    \item Zabezpieczenie hasłem Każda jednostka powinna być zabezpieczona hasłem charakterystycznym dla siebie, innym od domyślnego hasła. Hasła powinny spełniać poniższe wymogi:
    \begin{itemize}
        \item zawierać przynajmniej jedną wielkiej litery
        \item zawierać przynajmniej jedną małą literę
        \item zaiwerać przynajmniej jeden znak specjalny
        \item być nie krótsze niż 15 znaków
    \end{itemize}
    Każde hasło powinno również być zmieniane nie rzadziej niż co kwartał, w celu wyeliminowania możliwości długoterminowego dostępu do urządzeń w przypadku niewykrytego wycieku jednego lub kilku haseł.
    \item Wyłączenie nieużywanych funkcji dostępowych Wszystkie nieużywane funkcje dostępowe, takie jak zdalny dostęp przez WAN powinny zostać wyłączone. Wyeliminuje to potencjalne rejony ataków, ograniczając tym samym zakres ustawień oraz działań zabezpieczeń sieciowych.
    \item Regularne aktualizacje firmware'u Aby zapobiegać znanym błędom firmware'u zaleca się regularne jego aktualizacje. Można to zrealizować poprzez automatyczne aktualizacje, lub w celu zapewnienia stabilności działania infrastruktury, sporządzenie harmonogramu manualnych aktualizacji. Aktualizacje powinny być pobierane tylko z oficjalnych źródeł, takich jak strona internetowa producenta urządzenia.
    \item Dokumentacja konfiguracji sieciowej Podczas konfigurowania sieci równolegle powinna zostać sporządzona jej dokumentacja. W dokumentacji powinny zawierać się szczegóły konfiguracji, takie jak topologia sieci, hasła oraz wykorzystane urządzenia. Dokumentacja powinna być przetrzymywana lokalnie w zaszyfrowanej postaci i być udostępniana tylko upoważnionemu personelowi.
\end{enumerate}

\subsubsection{Drukarki i urządzenia
wielofunkcyjne}\label{drukarki-i-urzux105dzenia-wielofunkcyjne}

Drukarki oraz urządzenia wielofunkcyjne również zostały objęte procedurami bezpieczeństwa. Urządzenia te często działają jak tymczasowy bufor przechowujący dane wrażliwe, przez co mogą się stać potencjalnym celem ataków. Aby temu zapobiec sporządzono poniższą listę wymogów bezpieczeństwa:
\begin{enumerate}
    \item Zabezpieczenie hasłem Każde urządzenie powinno zostać zabezpieczone hasłem charakterystycznym dla siebie, innym od domyślnego hasła. Hasła powinny spełniać poniższe wymogi:
    \begin{itemize}
        \item zawierać przynajmniej jedną wielkiej litery
        \item zawierać przynajmniej jedną małą literę
        \item zawierać przynajmniej jeden znak specjalny
        \item być nie krótsze niż 15 znaków
    \end{itemize}
    \item Wyłączenie nieużywanych funkcji dostępowych Wszystkie nieużywane funkcje dostępowe, takie jak FTP powinny zostać wyłączone. Wyeliminuje to potencjalne rejony ataków, ograniczając tym samym zakres ustawień oraz działań zabezpieczeń sieciowych.
    \item Szyfrowanie połączeń W przypadku wykorzystania funkcjonalności sieciowych, połączenia z urządzeniami powinny być szyfrowane, niezależnie od funkcjonalności danej jednostki.
    \item Włączenie bezpiecznego wydruku Włączenie bezpiecznego wydruku dodatkowo będzie wymagać kodu PIN przed wykonaniem każdej operacji. Urządzenia powinny posiadać kod PIN charakterystyczny dla siebie. Kod do danego urządzenia powinien być udostępniony tylko osobom upoważnionym do jego obsługi.
    \item Regularne czyszczenie pamięci wewnętrznej Urządzenia, które mają możliwość cache'owania przekazywanej im zawartości, takie jak drukarki, powinny być regularnie czyszczone. Niedopuszczalnym jest aby pamięć urządznia pozostawała zapełniona dłużej niż przez pięć dni roboczych.

\end{enumerate}

\subsection{Harmonogram weryfikacji -- lista
kontrolna}\label{harmonogram-weryfikacji-lista-kontrolna}

\begin{longtable}[]{@{}>{\raggedright\arraybackslash}p{\dimexpr\linewidth/1000*270-2\tabcolsep\relax}>{\raggedright\arraybackslash}p{\dimexpr\linewidth/1000*430-2\tabcolsep\relax}>{\raggedright\arraybackslash}p{\dimexpr\linewidth/1000*300-2\tabcolsep\relax}@{}}
\toprule\noalign{}
Czynność weryfikacyjna & Częstotliwość & Status {[} {]} / Data \\
\midrule\noalign{}
\endhead
\bottomrule\noalign{}
\endlastfoot
Aktualizacja firmware routerów i switchy & Co 3 miesiące & {[} {]}
\underline{\hspace{2cm}} \\
Weryfikacja aktywnego oprogramowania antywirusowego & Co miesiąc & {[}
{]} \underline{\hspace{2cm}} \\
Skan podatności sieci & Co 6 miesięcy & {[} {]}
\underline{\hspace{2cm}} \\
Przegląd listy urządzeń podłączonych do sieci & Co miesiąc & {[} {]}
\underline{\hspace{2cm}} \\
Weryfikacja szyfrowania dysków (BitLocker/FileVault) & Co 6 miesięcy &
{[} {]} \underline{\hspace{2cm}} \\
Aktualizacja inwentarza sprzętowego & Co 12 miesięcy & {[} {]}
\underline{\hspace{2cm}} \\
Przegląd konfiguracji firewall & Co 6 miesięcy & {[} {]}
\underline{\hspace{2cm}} \\
\end{longtable}

\subsection{Rejestr sprzętu
(szablon)}\label{rejestr-sprzux119tu-szablon}

\begin{longtable}[]{@{}>{\raggedright\arraybackslash}p{\dimexpr\linewidth/1000*180-2\tabcolsep\relax}>{\raggedright\arraybackslash}p{\dimexpr\linewidth/1000*170-2\tabcolsep\relax}>{\raggedright\arraybackslash}p{\dimexpr\linewidth/1000*170-2\tabcolsep\relax}>{\raggedright\arraybackslash}p{\dimexpr\linewidth/1000*170-2\tabcolsep\relax}>{\raggedright\arraybackslash}p{\dimexpr\linewidth/1000*170-2\tabcolsep\relax}>{\raggedright\arraybackslash}p{\dimexpr\linewidth/1000*140-2\tabcolsep\relax}@{}}
\toprule\noalign{}
Urządzenie & Nr seryjny & Użytkownik & Data zakupu & Ostatni przegląd &
Status \\
\midrule\noalign{}
\endhead
\bottomrule\noalign{}
\endlastfoot
\underline{\hspace{2cm}} & \underline{\hspace{2cm}} &
\underline{\hspace{2cm}} & \underline{\hspace{2cm}} &
\underline{\hspace{2cm}} & {[} {]} OK \\
\underline{\hspace{2cm}} & \underline{\hspace{2cm}} &
\underline{\hspace{2cm}} & \underline{\hspace{2cm}} &
\underline{\hspace{2cm}} & {[} {]} OK \\
\underline{\hspace{2cm}} & \underline{\hspace{2cm}} &
\underline{\hspace{2cm}} & \underline{\hspace{2cm}} &
\underline{\hspace{2cm}} & {[} {]} OK \\
\end{longtable}

\begin{longtable}[]{@{}
  >{\raggedright\arraybackslash}p{\dimexpr\linewidth/1000*500-1\tabcolsep\relax}
  >{\raggedright\arraybackslash}p{\dimexpr\linewidth/1000*500-1\tabcolsep\relax}@{}}
\toprule\noalign{}
\endhead
\bottomrule\noalign{}
\endlastfoot
\begin{minipage}[t]{\linewidth}\raggedright
Zatwierdził:

Podpis: \underline{\hspace{3cm}}

Data: \underline{\hspace{4cm}}
\end{minipage} & \begin{minipage}[t]{\linewidth}\raggedright
Stanowisko / Tytuł:

Imię i nazwisko: \underline{\hspace{3cm}}

Funkcja: \underline{\hspace{3cm}}
\end{minipage} \\
\end{longtable}
\newpage
\phantomsection\label{doc3}
PROC-006

\section{PROCEDURA WERYFIKACJI CERTYFIKACJI I BEZPIECZEŃSTWA
OPROGRAMOWANIA PROC-006}\label{procedura-weryfikacji-certyfikacji-i-bezpieczeux144stwa-oprogramowania}

\begin{longtable}[]{@{}>{\raggedright\arraybackslash}p{\dimexpr\linewidth/1000*300-2\tabcolsep\relax}>{\raggedright\arraybackslash}p{\dimexpr\linewidth/1000*700-2\tabcolsep\relax}@{}}
\toprule\noalign{}
\endhead
\bottomrule\noalign{}
\endlastfoot
Wersja: & 1.0 \\
Dotyczy: & Systemy operacyjne, aplikacje biznesowe, oprogramowanie
księgowe, przeglądarki \\
Odpowiedzialny: & Administrator IT \\
Częstotliwość: & Przy instalacji + kwartalny przegląd \\
\end{longtable}

\subsection{Lista Zatwierdzonego Oprogramowania
(LZO)}\label{lista-zatwierdzonego-oprogramowania-lzo}

W firmie dopuszczone jest wyłącznie oprogramowanie z LZO. Instalacja
oprogramowania spoza listy wymaga pisemnej zgody administratora IT.

\begin{longtable}[]{@{}>{\raggedright\arraybackslash}p{\dimexpr\linewidth/1000*270-2\tabcolsep\relax}>{\raggedright\arraybackslash}p{\dimexpr\linewidth/1000*250-2\tabcolsep\relax}>{\raggedright\arraybackslash}p{\dimexpr\linewidth/1000*250-2\tabcolsep\relax}>{\raggedright\arraybackslash}p{\dimexpr\linewidth/1000*230-2\tabcolsep\relax}@{}}
\toprule\noalign{}
Kategoria & Zatwierdzone rozwiązanie & Cel & Licencja OK \\
\midrule\noalign{}
\endhead
\bottomrule\noalign{}
\endlastfoot
System operacyjny & Windows 11 Pro / macOS 13+ & Stacje robocze & {[}
{]} TAK \\
Oprogramowanie księgowe & {[}NAZWA -- np. Insert/Symfonia{]} & Główna
aplikacja & {[} {]} TAK \\
Antywirus/EDR & {[}NAZWA{]} & Ochrona endpoint & {[} {]} TAK \\
VPN & {[}NAZWA{]} & Zdalny dostęp & {[} {]} TAK \\
Przeglądarka & Chrome / Firefox (aktualna wersja) & Internet & {[} {]}
TAK \\
Office & Microsoft 365 Business & Praca biurowa & {[} {]} TAK \\
Manager haseł & {[}NAZWA -- np. Bitwarden/1Password{]} & Zarządzanie
hasłami & {[} {]} TAK \\
\end{longtable}

\subsection{Procedura instalacji nowego
oprogramowania}\label{procedura-instalacji-nowego-oprogramowania}
Instalacja nowego, niezweryfikowanego oprogramowania wiąże się z potencjalnym ryzykiem ataku lub wycieku danych spowodowanym niewystarczającymi środkami bezpieczeństwa instalowanego oprogramowania. Pełne wyeliminowanie ryzyka nie jest możliwe, ale podjęcie odpowiednich kroków w celu jego zminimalizowania jest jak najbardziej zalecane.
\begin{enumerate}
\tightlist
\item
  \textbf{Wniosek pracownika do administratora IT (forma pisemna/email)}
  Wybrane oprogramowanie może zostać zgłoszone do administratora IT w celu jego weryfikacji. Administratorowi przysługuje wówczas 5 dni roboczych na zweryfikowanie i zatwierdzenie lub odrzucenie prośby. Aby zgłosić dany program do weryfikacji należy posłużyć się poniższym wzorem wniosku:

\begin{verbatim}
Chciałbym/ałabym zgłosić program o nazwie .................... 
do weryfikacji przez administratora IT,
w celu jego zainstalowania na komputerze 
o numerze identyfikacyjnym ..................... 
Program służy do .....................
i będzie wykorzystywany w celu ......................
\end{verbatim}

  Wniosek powinien zostać wysłany na wewnętrzny adres mailowy wsparcia technicznego organizacji. Odpowiedź administratora, wraz z uzasadnieniem (szczególnie w przypadku odmowy) powinna zostać odesłana na adres mailowy nadawcy. Dodatkowo powinna ona zostać zarchiwizowana na wypadek ponownej prośby w przyszłości.
\item
  \textbf{Ocena bezpieczeństwa}
  \newline
  Ocena bezpieczeństwa aplikacji powinna zostać przeprowadzona przez administratora IT lub upoważniony do tego personel za jego zgodą. Podczas oceny pod uwagę powinny być wzięte takie aspekty jak:
  \begin{enumerate}
      \item Pochodzenie Jeśli aplikacja pochodzi od nieznanego wydawcy, a jej szczegółowe działanie jest niemożliwe do przeanalizowania to powinna zostać odrzucona.
      \item Funkcjonalność aplikacji Jeśli aplikacja nie oferuje zadowalającej funkcjonalności lub jej funkcjonalność pokrywa się z programem, który jest już wpisany do LZO, aplikacja powinna zostać odrzucona.
      \item Stopień ryzyka Jeśli aplikacja zbiera dane użytkowników, lub jej funkcjonalność opiera się o połączenie internetowe (zapis w chmurze, wykorzystanie zewnętrznych serwerów) to powinna być szczegółowo przeanalizowana pod kątem ryzyka, które wprowadza do organizacji. Jeśli stopień ryzyka jest zbyt duży, to aplikacja powinna zostać odrzucona.
      \item Sprawdzenie CVE Punkt ten jest rozwinięciem poprzedniego punktu. CVE (Common Vulnerabilities and Exposures) pozwala zaznajomić się z listą powszechnie znanych podatności i zagrożeń. Jeśli lista informuje o zbyt dużym stopniu ryzyka to aplikacja powinna zostać odrzucona.
      \item Wsparcie Jeśli aplikacja nie jest udokumentowana ani nie posiada społeczności wsparcia to powinna zostać odrzucona.
  \end{enumerate}
\item
  \textbf{Weryfikacja licencji}
  \newline
  Przed przystąpieniem do instalacji oprogramowania powinien zostać zweryfikowany typ licencji na bazie której ono operuje. W przypadku licencji uniemożliwiającej instalację i swobodne użytkowanie w ramach działalności orgranizacji dalsze kroki powinny zostać zaniechane. W przypadku braku wykupionej płatnej licencji komercyjnej na dane oprogramowanie lub jego wersję, prośba o instalację powinna zostać zgłoszona do odpowiednich działów organizacji w celu jej zakupienia.
\item
  \textbf{Instalacja} 
  \newline
  Instalacja oprogramowania powinna zostać przeprowadzona przez administratora IT lub upoważniony do tego personel za jego zgodą. Dopuszcza się przeprowadzenie instalacji przez system zdalnego dostępu (TeamViewer, AnyDesk). Plik instalacyjny należy pobrać z oficjalnego źródła, takiego jak strona producenta, a procedurę instalacji przeprowadzić zgodnie z instrukcją obsługi instalowanego oprogramowania. Niedopuszczalnym jest aby użytkownik wypożyczonego sprzętu przeprowadził instalację oprogramowania we własnym zakresie.
\item
  \textbf{Dodanie do LZO i dokumentacja}
  \newline
  Po przeprowadzonej instalacji program powinien zostać dodany do LZO wraz z dokumentacją podjętych działań. Dokumentacja powinna zawierać raport z procesu instalacji oraz uzasadnienie decyzji o dopisaniu programu od LZO.
\item
  \textbf{Monitorowanie przez 30 dni po instalacji}
  \newline
  Każda nowa aplikacja powinna być monitorowana przez 30 dni po jej instalacji. Aspekty wymagające szczególnej uwagi powinny zawierać:
  \begin{enumerate}
      \item Stabilność działania stacji roboczej Aplikacja nie powinna destabilizować działania stacji roboczej na której została zainstalowana. Zdarzenia takie jak samoistne wyłączanie się komputera, spowolnione działanie czy zablokowanie niektórych funkcji lub programów powinny zostać bezzwłocznie zgłoszone do administratora IT.
      \item Zużycie zasobów stacji roboczej Zużycie zasobów stacji roboczej nie powinno wzrosnąć po instalacji aplikacji. Jeśli program zużywa nadmierne ilości zasobów powinno to zostać bezzwłocznie zgłoszone do administratora IT.
      \item Integralność danych na dysku twardym Dane oraz ich integralność na dysku twardym powinny być regularnie kontrolowane. Ma to na celu szybką reakcję w przypadku niekontrolowanej enkrypcji danych przez aplikację. Jakiekolwiek problemy z integralnością danych powinny zostać bezzwłocznie zgłoszone do administratora IT.
  \end{enumerate}
\end{enumerate}

\subsection{Aktualizacje i łatki
bezpieczeństwa}\label{aktualizacje-i-ux142atki-bezpieczeux144stwa}
Regularne aktualizacje aplikacji oraz systemów operacyjnych są równie ważne dla bezpieczeństwa organizacji, co aktualizacje firmware'u kluczowego osprzętu infrastruktury sieciowej. Proces oraz harmonogram aktualizacji oprogramowania w obrębie orgranizacji powinien być podyktowany poniższymi zasadami:
\begin{enumerate}
    \item \textbf{Aktualizacje systemu operacyjnego}
    \newline
    Standardowe aktualizacje systemu mogą być przeprowadzane według harmonogramu dystrybutora oprogramowania lub manualnie, według harmonogramu aktualizacji orgranizacji. Aktualizacje krytyczne, zawierające łatki na znane, poważne luki bezpieczeństwa powinny być zainstalowane nie później niż 72h po ich wydaniu. Wszystkie aktualizacje powinny pochodzić z oficjalnych źródeł, takich jak strona internetowa producenta systemu operacyjnego.
    \item \textbf{Oprogramowanie biznesowe}
    Aktualizacje oprogramowania biznesowego powinny być rozpatrywane indywidualnie, z uwzględnieniem funkcjonalności oprogramowania oraz zmian jakie wprowadza dana łatka. Aktualizacje niezwiązane z bezpieczeństwem, takie jak zmiany interfejsu graficznego lub innych funkcji aplikacji mogą zostać pominięte. Aktualizacje krytyczne dla oprogramowania narażonego na ataki czy wycieki danych powinny być zainstalowane nie później niż 72h po ich wydaniu. Aktualizacje powinny pochodzić z oficjalnych źródeł, takich jak strona internetowa producenta oprogramowania.
    \item \textbf{Przeglądarki i wtyczki}
    Przeglądarki należą do grupy aplikacji na stale podłączonych do internetu i wykonujących cudze programy na maszynie odwiedzającego. Z tego powodu jakiekolwiek ich aktualizacje powinny być instalowane nie później niż 72h po ich wydaniu.
    \item \textbf{Oprogramowanie bez aktualizacji od >12 miesięcy}
    Oprogramowanie, które utraciło wsparcie powinno zostać odcięte od dostępu do jakiejkolwiek sieci. Jeśli aplikacja opiera się na funkcjach sieciowych powinna zostać bezzwłocznie odinstalowana. W przypadku silnej zależności firmy od takiej aplikacji, administrator IT powinien uważnie monitorować działanie komputerów na których jest ona zainstalowana. Organizacja powinna podjąć kroki aby jak najszybciej znaleźć oraz wdrożyć alternatywę dla oprogramowania bez wsparcia.
\end{enumerate}


\subsection{4. Lista kontrolna -- kwartalny przegląd
oprogramowania}\label{lista-kontrolna-kwartalny-przeglux105d-oprogramowania}

\begin{longtable}[]{@{}>{\raggedright\arraybackslash}p{\dimexpr\linewidth/1000*270-2\tabcolsep\relax}>{\raggedright\arraybackslash}p{\dimexpr\linewidth/1000*430-2\tabcolsep\relax}>{\raggedright\arraybackslash}p{\dimexpr\linewidth/1000*300-2\tabcolsep\relax}@{}}
\toprule\noalign{}
Czynność & Termin & Status \\
\midrule\noalign{}
\endhead
\bottomrule\noalign{}
\endlastfoot
Weryfikacja wersji systemu operacyjnego na wszystkich komputerach & Co
kwartał & {[} {]} \underline{\hspace{2cm}} \\
Przegląd listy zainstalowanego oprogramowania (porównanie z LZO) & Co
kwartał & {[} {]} \underline{\hspace{2cm}} \\
Sprawdzenie dat wygaśnięcia licencji & Co kwartał & {[} {]}
\underline{\hspace{2cm}} \\
Weryfikacja bazy sygnatur antywirusa & Co miesiąc & {[} {]}
\underline{\hspace{2cm}} \\
Skan podatności zainstalowanego oprogramowania (CVE) & Co 6 miesięcy &
{[} {]} \underline{\hspace{2cm}} \\
Usunięcie nieautoryzowanego lub nieaktualnego SW & Bieżąco & {[} {]}
\underline{\hspace{2cm}} \\
\end{longtable}

\begin{longtable}[]{@{}
  >{\raggedright\arraybackslash}p{\dimexpr\linewidth/1000*500-1\tabcolsep\relax}
  >{\raggedright\arraybackslash}p{\dimexpr\linewidth/1000*500-1\tabcolsep\relax}@{}}
\toprule\noalign{}
\endhead
\bottomrule\noalign{}
\endlastfoot
\begin{minipage}[t]{\linewidth}\raggedright
Zatwierdził:

Podpis: \underline{\hspace{3cm}}

Data: \underline{\hspace{4cm}}
\end{minipage} & \begin{minipage}[t]{\linewidth}\raggedright
Stanowisko / Tytuł:

Imię i nazwisko: \underline{\hspace{3cm}}

Funkcja: \underline{\hspace{3cm}}
\end{minipage} \\
\end{longtable}

\phantomsection\label{doc4}
PROC-007

\section{PROCEDURA ZABEZPIECZANIA INFRASTRUKTURY PRZED
ATAKAMI PROC-007}\label{procedura-zabezpieczania-infrastruktury-przed-atakami}

\begin{longtable}{@{}>{\raggedright\arraybackslash}p{\dimexpr\linewidth/1000*300-2\tabcolsep\relax}>{\raggedright\arraybackslash}p{\dimexpr\linewidth/1000*700-2\tabcolsep\relax}@{}}
\toprule\noalign{}
\endhead
\bottomrule\noalign{}
\endlastfoot
Wersja: & 1.0 \\
Dotyczy: & Sieć firmowa, serwery, stacje robocze, dostęp zdalny \\
Odpowiedzialny: & Administrator IT \\
Przegląd: & Co 6 miesięcy lub po każdym istotnym incydencie bezpieczeństwa \\
Norma odniesienia: & ISO/IEC 27001:2022, Załącznik A kontrole: A.8.20, A.8.21, A.8.22, A.8.23, A.8.9, A.8.5, A.8.7, A.8.8, A.8.13, A.8.15, A.7.1, A.7.2 \\
\end{longtable}

\subsection{Cel i zakres procedury}\label{cel-i-zakres-procedury}

Niniejsza procedura określa wymagania dotyczące zabezpieczenia infrastruktury teleinformatycznej firmy przed zagrożeniami zewnętrznymi i wewnętrznymi. Procedura realizuje wymagania normy ISO/IEC 27001:2022 w zakresie kontroli technologicznych (Załącznik A, kategoria 8) oraz jest zgodna z wymaganiami ustawy o krajowym systemie cyberbezpieczeństwa. Zakres obejmuje: sieć lokalną (LAN/WLAN), urządzenia sieciowe (router, firewall, switch), serwery i stacje robocze, systemy dostępu zdalnego (VPN), systemy kopii zapasowych, oprogramowanie ochronne (EDR/AV/DNS filtering).



\subsection{Architektura sieciowa}\label{architektura-sieciowa}

Zgodnie z kontrolą \textbf{A.8.22 (Segregacja sieci)} sieć firmy jest podzielona na odseparowane segmenty:
\begin{itemize}
\tightlist
\item
  \textbf{Sieć biurowa (LAN):} stacje robocze pracowników -- VLAN 10.
\item
  \textbf{Sieć serwerowa:} serwery plików, baz danych i aplikacji księgowych -- VLAN 20, dostęp wyłącznie z LAN i VPN.
\item
  \textbf{Sieć dla gości (Guest WiFi):} odizolowana od sieci wewnętrznej, brak dostępu do zasobów firmowych -- VLAN 30.
\item
  \textbf{Strefa DMZ:} dla usług wystawionych na zewnątrz.
\end{itemize}

Firewall na granicy sieci stosuje politykę domyślnej odmowy (default deny) -- cały ruch jest blokowany, chyba że jawnie dozwolony przez regułę. Brak bezpośredniego dostępu RDP/SSH z Internetu -- wymagany VPN. Dokumentacja topologii sieci jest aktualizowana przy każdej zmianie konfiguracji i przechowywana w sposób poufny (wydruk + kopia zaszyfrowana).

\subsection{Zarządzanie podatnościami technicznymi}\label{zarzadzanie-podatnosciami-technicznymi}

Zgodnie z kontrolą \textbf{A.8.8 (Zarządzanie podatnościami technicznymi)}:
\begin{itemize}
\tightlist
\item
  Skany podatności sieci wewnętrznej i zewnętrznej: co najmniej co 6 miesięcy (narzędzia: OpenVAS, Nessus lub równoważne).
\item
  Łatki krytyczne (CVSS $\geq$ 9.0): wdrożenie w ciągu 48 godzin.
\item
  Łatki wysokie (CVSS 7.0--8.9): wdrożenie w ciągu 7 dni.
\item
  Łatki średnie (CVSS 4.0--6.9): wdrożenie w ciągu 30 dni.
\item
  Łatki niskie (CVSS $<$ 4.0): wdrożenie przy najbliższym oknie serwisowym.
\item
  Rejestr wdrożonych łatek prowadzony przez administratora IT.
\item
  Oprogramowanie bez wsparcia producenta (EOL): plan migracji w ciągu 90 dni lub akceptacja ryzyka przez właściciela firmy.
\end{itemize}

\subsection{Hasła i
uwierzytelnianie}\label{hasux142a-i-uwierzytelnianie}

Zgodnie z kontrolą \textbf{A.8.5 (Bezpieczne uwierzytelnianie)}:

\begin{itemize}
\tightlist
\item
  Minimalna długość hasła: \textbf{12 znaków} (konta standardowe), \textbf{16 znaków} (konta uprzywilejowane)
\item
  Złożoność: wielkie/małe litery, cyfry, znaki specjalne
\item
  Zmiana haseł: co 45 dni (konta uprzywilejowane) / co 90 dni
  (standardowe)
\item
  Obowiązkowe \textbf{MFA} dla: kont administratora, VPN, poczty e-mail,
  systemów księgowych, dostępu do kopii zapasowych
\item
  Zarządzanie hasłami przez zatwierdzonego managera haseł (np. Bitwarden,1Password, KeePassXC)
\item
  Blokada konta po \textbf{5 błędnych próbach} logowania
\item 
  Zakaz stosowania tych samych hasel w systemach firmowych i prywatnych
\item 
  Zakaz przechowywania haseł w formie jawnej (pliki tekstowe, karteczki, arkusze
kalkulacyjne)

\end{itemize}

\subsection{Dostęp zdalny}\label{dostux119p-zdalny}

Zgodnie z kontrolą \textbf{A.8.20 (Bezpieczeństwo sieci)} i \textbf{A.8.21 (Bezpieczeństwo usług sieciowych)}:

\begin{itemize}
\tightlist
\item
  Zdalny dostęp wyłącznie przez VPN (z MFA)
\item
  Lista zatwierdzonych użytkowników VPN -- przegląd co 3 miesiące
\item
  Sesje VPN z limitem czasu (timeout po 30 min bezczynności)
\item
  Zakaz używania publicznych sieci WiFi
\item 
  Logowanie wszystkich sesji VPN: data, godzina, adres IP zródłowy, uzytkownik,czas trwania
\item 
  Zatwierdzenie dostępu zdalnego wymaga pisemnego wniosku zaakceptowanego przez właściciela firmy

\end{itemize}

\subsection{Ochrona przed złośliwym
oprogramowaniem}\label{Ochrona-przed-złośliwym-oprogramowaniem}

Zgodnie z kontrolą \textbf{A.8.7 (Ochrona przed złośliwym oprogramowaniem)}:

\begin{itemize}
\tightlist
\item
  Oprogramowanie EDR/Antywirus na każdym urządzeniu końcowym
\item
  Blokowanie wykonywania skryptów z makrami w dokumentach Office
  (polityka grupowa)
\item
  Filtrowanie DNS z blokowaniem złośliwych domen
\item
  Polityka białej listy aplikacji (Application Whitelisting) na stacjach
  roboczych
\item
  Skanowanie załączników e-mail przez antywirus
\item 
  Zakaz uruchamiania plików wykonywalnych z katalogów tymczasowych i pobierania

\end{itemize}

\subsection{Polityka kopii zapasowych
(3-2-1)}\label{polityka-kopii-zapasowych-3-2-1}

Zgodnie z kontrolą \textbf{A.8.13 (Kopie zapasowe informacji)}:

Reguła 3-2-1: \textbf{3 kopie}, na \textbf{2 różnych nośnikach}, w tym
\textbf{1 offsite}. Testowanie przywracania: co kwartał.

\begin{longtable}[]{@{}>{\raggedright\arraybackslash}p{\dimexpr\linewidth/1000*270-2\tabcolsep\relax}>{\raggedright\arraybackslash}p{\dimexpr\linewidth/1000*250-2\tabcolsep\relax}>{\raggedright\arraybackslash}p{\dimexpr\linewidth/1000*250-2\tabcolsep\relax}>{\raggedright\arraybackslash}p{\dimexpr\linewidth/1000*230-2\tabcolsep\relax}@{}}
\toprule\noalign{}
Rodzaj backup & Częstotliwość & Przechowywanie & Szyfrowanie \\
\midrule\noalign{}
\endhead
\bottomrule\noalign{}
\endlastfoot
Pełny backup serwerów & Co tydzień (niedziela) & 90 dni & TAK \\
Przyrostowy backup plików & W czasie rzeczywistym & 30 dni & TAK \\
Backup baz danych (księgowość) & Co 4 godziny & 90 dni &
TAK \\
Backup poza siedzibą (offsite) & Co tydzień & 12 miesięcy &
TAK \\
Backup konfiguracji sieci & Po każdej zmianie & Bezterminowo & TAK \\
\end{longtable}

\subsection{Bezpieczeństwo fizyczne infrastruktury}\label{Bezpieczeństwo-fizyczne-infrastruktury}

Zgodnie z kontrolami \textbf{A.7.1 (Fizyczne granice bezpieczeństwa)} i \textbf{A.7.2 (Fizyczne kontrole wejścia)}:

\begin{itemize}
\tightlist
\item
  Serwery i urządzenia sieciowe umieszczone w zamykanej szafie serwerowej lub wydzielonym pomieszczeniu z kontrolą dostępu
\item
  Klucze/karty dostępu do pomieszczenia serwerowni: wyłączenie administrator IT i właściciel firmy
\item
  Rejestr wejść do pomieszczenia serwerowni
\item
  Zabezpieczenie przed czynnikami środowiskowymi: UPS, generator prądu, ochrona przeciwprzepięciowa, odpowiednia wentylacja/klimatyzacja, ochrona przeciwpożarowa
\item 
  Zakaz przechowywania materiałów łatwopalnych w pobliżu sprzętu IT

\end{itemize}

\subsection{Zarządzanie zmianami w infrastrukturze}\label{Zarządzanie-zmianami-w-infrastrukturze}

Zgodnie z kontrolą \textbf{A.8.9 (Zarządzanie konfiguracją)}:

\begin{itemize}
\tightlist
\item
 Każda zmiana w konfiguracji sieci, serwerów lub systemów bezpieczeństwa wymaga udokumentowanego wniosku o zmianę (Change Request)
\item
  Wniosek musi zawierać opis zmiany, uzasadnienie, ocenę wpływu na bezpieczenstwo,
plan wycofania (rollback plan)
\item
  Zatwierdzenie zmiany: administrator IT + właściciel firmy (dla zmian krytycznych)
\item
 Dokumentacja wdrożonej zmiany: data, osoba odpowiedzialna, opis techniczny, wynik testów
\item 
  Bazowa konfguracja (baseline) utrzymywana i aktualizowana po każdej zatwierdzonej zmianie

\end{itemize}

\subsection{Rejestr zmian i przeglądów procedury}\label{rejestr-zmian-i-przegladow-procedury}

\begin{longtable}{@{}>{\raggedright\arraybackslash}p{\dimexpr\linewidth/1000*150-2\tabcolsep\relax}>{\raggedright\arraybackslash}p{\dimexpr\linewidth/1000*200-2\tabcolsep\relax}>{\raggedright\arraybackslash}p{\dimexpr\linewidth/1000*450-2\tabcolsep\relax}>{\raggedright\arraybackslash}p{\dimexpr\linewidth/1000*200-2\tabcolsep\relax}@{}}
\toprule\noalign{}
Wersja & Data & Opis zmiany & Autor \\
\midrule\noalign{}
\endhead
\bottomrule\noalign{}
\endlastfoot
1.0 & {} & Wersja początkowa (Dostosowanie do ISO/IEC 27001:2022) & {} \\
\end{longtable}

\begin{longtable}{@{}
  >{\raggedright\arraybackslash}p{\dimexpr\linewidth/1000*500-1\tabcolsep\relax}
  >{\raggedright\arraybackslash}p{\dimexpr\linewidth/1000*500-1\tabcolsep\relax}@{}}
\toprule\noalign{}
\endhead
\bottomrule\noalign{}
\endlastfoot
\begin{minipage}[t]{\linewidth}\raggedright
Zatwierdził:

Podpis: \underline{\hspace{3cm}}

Data: \underline{\hspace{4cm}}
\end{minipage} & \begin{minipage}[t]{\linewidth}\raggedright
Stanowisko / Tytuł:

Imię i nazwisko: \underline{\hspace{3cm}}

Funkcja: \underline{\hspace{3cm}}
\end{minipage} \\
\end{longtable}

\phantomsection\label{doc5}
POL-008


\section{POLITYKA MONITOROWANIA I AUDYTOWANIA
INFRASTRUKTURY POL-008}\label{polityka-monitorowania-i-audytowania-infrastruktury}

\begin{longtable}{@{}>{\raggedright\arraybackslash}p{\dimexpr\linewidth/1000*300-2\tabcolsep\relax}>{\raggedright\arraybackslash}p{\dimexpr\linewidth/1000*700-2\tabcolsep\relax}@{}}
\toprule\noalign{}
\endhead
\bottomrule\noalign{}
\endlastfoot
Wersja: & 1.0 \\
Dotyczy: & Wszystkie systemy IT firmy, urządzenia sieciowe, usługi chmurowe \\
Odpowiedzialny: & Administrator IT \\
Przegląd: & Co 12 miesięcy lub po istotnym incydencie bezpieczeństwa \\
Norma odniesienia: & ISO/IEC 27001:2022, Załącznik A kontrole: A.8.15, A.8.16, A.8.17, A.5.35, A.5.36, A.8.34 \\
\end{longtable}

\subsection{Cel i zakres polityki}\label{Cel-i-zakres-polityki}

Niniejsza polityka określa zasady monitorowania, rejestrowania zdarzeń (logowania) oraz audytowania infrastruktury IT firmy. Celem jest zapewnienie ciągłego nadzoru nad bezpieczeństwem systemów informatycznych, wykrywanie anomalii i naruszeń oraz spełnienie wymagań normy ISO/IEC 27001:2022 w zakresie kontroli \textbf{A.8.15 (Logowanie)}, \textbf{A.8.16 (Działania monitorujące)} i \textbf{A.8.17 (Synchronizacja zegarów)}. Polityka obejmuje: serwery, stacje robocze, urządzenia sieciowe (firewall, router, switch), systemy pocztowe, systemy księgowe, usługi VPN, usługi chmurowe.



\subsection{Logi systemowe - zakres zbieranych danych}\label{logi-systemowe-zakres-zbieranych-danych}

Zgodnie z kontrolą \textbf{A.8.15 (Logowanie)} rejestrowane są co najmniej następujące zdarzenia:

\begin{longtable}{@{}>{\raggedright\arraybackslash}p{\dimexpr\linewidth/1000*200-2\tabcolsep\relax}>{\raggedright\arraybackslash}p{\dimexpr\linewidth/1000*350-2\tabcolsep\relax}>{\raggedright\arraybackslash}p{\dimexpr\linewidth/1000*150-2\tabcolsep\relax}>{\raggedright\arraybackslash}p{\dimexpr\linewidth/1000*150-2\tabcolsep\relax}>{\raggedright\arraybackslash}p{\dimexpr\linewidth/1000*150-2\tabcolsep\relax}@{}}
\toprule\noalign{}
Źródło logów & Rodzaj rejestrowanych zdarzeń & Retencja & Szyfrow. & Kontrola ISO \\
\midrule\noalign{}
\endhead
\bottomrule\noalign{}
\endlastfoot
Serwer Windows/Linux & Logowania (sukces/porażka), zmiany uprawnień, zdarzenia systemowe, instalacje oprogramowania & 12 mies. & TAK & A.8.15 \\
Firewall / Router & Połączenia przychodzące i wychodzące, blokady reguł, anomalie ruchu, zmiany konfiguracji & 12 mies. & TAK & A.8.15, A.8.20 \\
Poczta e-mail & Wysyłane/odebrane wiadomości (metadane), spam, podejrzane załączniki, próby phishingu & 12 mies. & TAK & A.8.15 \\
System księgowy & Logowania, zmiany w danych klientów, generowanie raportów, eksport danych & 5 lat & TAK & A.8.15, A.5.33 \\
Stacje robocze (EDR) & Procesy, połączenia sieciowe, podłączenia USB, uruchomienia aplikacji & 90 dni & TAK & A.8.7, A.8.15 \\
VPN & Sesje (start/stop), adres IP źródłowy, użytkownik, czas trwania, wolumen danych & 12 mies. & TAK & A.8.15, A.8.20 \\
Kontrola dostępu fizycznego & Wejścia/wyjścia do serwerowni, zdarzenia alarmowe & 12 mies. & TAK & A.7.2, A.8.15 \\
\end{longtable}

Każdy rekord logu musi zawierać co najmniej: datę i czas zdarzenia (zsynchronizowany), identyfikator użytkownika lub systemu, rodzaj zdarzenia, wynik zdarzenia (sukces/porażka), adres IP źródłowy.

\subsection{Synchronizacja zegarów}\label{synchronizacja-zegarow}

Zgodnie z kontrolą \textbf{A.8.17 (Synchronizacja zegarów)}:
\begin{itemize}
\tightlist
\item
  Wszystkie systemy i urządzenia synchronizowane z protokołem NTP (Network Time Protocol).
\item
  Źródło czasu: serwer NTP firmy lub publiczny serwer referencyjny (np. \texttt{tempus1.gum.gov.pl} -- Główny Urząd Miar).
\item
  Weryfikacja poprawności synchronizacji: co miesiąc (odchylenie dopuszczalne $<$ 1 sekunda).
\item
  Dokumentacja źródła czasu w rejestrze konfiguracji sieci.
\end{itemize}

\subsection{Monitorowanie w czasie rzeczywistym i alerty}\label{monitorowanie-w-czasie-rzeczywistym-i-alerty}

Zgodnie z kontrolą \textbf{A.8.16 (Działania monitorujące)}:

\subsubsection{Alerty automatyczne -- wyzwalacze (triggers)}\label{alerty-automatyczne-wyzwalacze-triggers}

\begin{longtable}{@{}>{\raggedright\arraybackslash}p{\dimexpr\linewidth/1000*400-2\tabcolsep\relax}>{\raggedright\arraybackslash}p{\dimexpr\linewidth/1000*400-2\tabcolsep\relax}>{\raggedright\arraybackslash}p{\dimexpr\linewidth/1000*200-2\tabcolsep\relax}@{}}
\toprule\noalign{}
Zdarzenie & Reakcja systemu & Priorytet \\
\midrule\noalign{}
\endhead
\bottomrule\noalign{}
\endlastfoot
Logowanie poza godzinami pracy (18:00--6:00 lub weekend) & Alert natychmiastowy do administratora IT & Wysoki \\
5+ nieudanych prób logowania w ciągu 10 min & Alert + tymczasowa blokada konta (30 min) & Wysoki \\
Dostęp do plików spoza normalnego zakresu pracy użytkownika & Alert do administratora IT & Średni \\
Wysyłka dużej liczby e-maili (\textgreater{}50/godz.) z jednego konta & Alert + tymczasowa blokada wysyłki & Wysoki \\
Transfer danych \textgreater{}1 GB w krótkim czasie & Alert do administratora IT & Wysoki \\
Nowe, niezidentyfikowane urządzenie w sieci (rogue device) & Alert natychmiastowy + izolacja portu (jeśli możliwe) & Krytyczny \\
Zmiana konfiguracji firewall/routera & Alert do administratora IT & Wysoki \\
Wykrycie złośliwego oprogramowania (EDR/Antywirus) & Alert natychmiastowy + automatyczna kwarantanna & Krytyczny \\
Nieudana kopia zapasowa & Alert do administratora IT & Wysoki \\
\end{longtable}

\subsubsection{Eskalacja alertów}\label{eskalacja-alertow}
\begin{itemize}
\tightlist
\item
  \textbf{Alerty Krytyczne:} natychmiastowe powiadomienie administratora IT i właściciela firmy (SMS/telefon); reakcja w ciągu 30 minut; uruchomienie procedury PLAY-010 w razie potrzeby.
\item
  \textbf{Alerty Wysokie:} powiadomienie administratora IT (e-mail + aplikacja); analiza i reakcja w ciągu 6 godzin.
\item
  \textbf{Alerty Średnie:} rejestracja w systemie; analiza przy najbliższym przeglądzie logów (do 24 godzin).
\end{itemize}

\subsection{Harmonogram przeglądów i audytów}\label{harmonogram-przegladow-i-audytow}

Zgodnie z kontrolami \textbf{A.5.35 (Niezależny przegląd bezpieczeństwa informacji)} i \textbf{A.5.36 (Zgodność z politykami i standardami)}:

\begin{longtable}{@{}>{\raggedright\arraybackslash}p{\dimexpr\linewidth/1000*300-2\tabcolsep\relax}>{\raggedright\arraybackslash}p{\dimexpr\linewidth/1000*150-2\tabcolsep\relax}>{\raggedright\arraybackslash}p{\dimexpr\linewidth/1000*200-2\tabcolsep\relax}>{\raggedright\arraybackslash}p{\dimexpr\linewidth/1000*150-2\tabcolsep\relax}>{\raggedright\arraybackslash}p{\dimexpr\linewidth/1000*200-2\tabcolsep\relax}@{}}
\toprule\noalign{}
Rodzaj audytu/przeglądu & Częstotl. & Odpowiedzialny & Kontrola ISO & Ostatni / Następny \\
\midrule\noalign{}
\endhead
\bottomrule\noalign{}
\endlastfoot
Przegląd logów bezpieczeństwa & Codziennie & Admin IT & A.8.15, A.8.16 & \underline{\hspace{1.5cm}} / \underline{\hspace{1.5cm}} \\
Weryfikacja aktywnych kont użytkowników & Co miesiąc & Admin IT & A.5.18 & \underline{\hspace{1.5cm}} / \underline{\hspace{1.5cm}} \\
Audyt uprawnień (porównanie z macierzą dostępu) & Co kwartał & Admin IT & A.5.15, A.8.3 & \underline{\hspace{1.5cm}} / \underline{\hspace{1.5cm}} \\
Skan podatności sieci (wewnętrzny /zewnętrzny) & Co kwartał & Admin IT / zewn. & A.8.8 & \underline{\hspace{1.5cm}} / \underline{\hspace{1.5cm}} \\
Przegląd konfiguracji firewall i reguł & Co 6 mies. & Admin IT & A.8.20, A.8.22 & \underline{\hspace{1.5cm}} / \underline{\hspace{1.5cm}} \\
Przegląd skuteczności kopii zapasowych & Co kwartał & Admin IT & A.8.13 & \underline{\hspace{1.5cm}} / \underline{\hspace{1.5cm}} \\
Test penetracyjny (pen-test) & Co 12 mies. & Zewnętrzny & A.8.8, A.5.35 & \underline{\hspace{1.5cm}} / \underline{\hspace{1.5cm}} \\
Audyt zgodności z RODO & Co 12 mies. & IOD / zewnętrzny & A.5.36 & \underline{\hspace{1.5cm}} / \underline{\hspace{1.5cm}} \\
Niezależny przegląd ISMS & Co 12 mies. & Zewnętrzny audytor & A.5.35 & \underline{\hspace{1.5cm}} / \underline{\hspace{1.5cm}} \\
Przegląd polityk i procedur bezpieczeństwa & Co 12 mies. & Właściciel / Admin IT & A.5.1 & \underline{\hspace{1.5cm}} / \underline{\hspace{1.5cm}} \\
\end{longtable}

\newpage
\subsection{Przechowywanie i ochrona logów}\label{przechowywanie-i-ochrona-logow}

Zgodnie z kontrolą \textbf{A.8.15 (Logowanie)} oraz wymogami rozliczalności z art. 5 ust. 2 RODO:
\begin{itemize}
\tightlist
\item
  Logi przechowywane na oddzielnym, zabezpieczonym serwerze.
\item
  Dostęp do logów: wyłącznie administrator IT i właściciel firmy (zasada minimalnych uprawnień).
\item
  Logi niemodyfikowalne: zapis w trybie write-once (WORM).
\item
  Kopia logów w chmurze (zaszyfrowana AES-256) przechowywana przez okres retencji.
\item
  Logi zawierające dane osobowe (np. identyfikatory użytkowników) podlegają ochronie zgodnie z RODO - dostęp ograniczony, pseudonimizacja tam, gdzie możliwa.
\item
  Usuwanie logów po upływie okresu retencji - udokumentowane, z zachowaniem procedury bezpiecznego usuwania danych.
\end{itemize}

\subsection{Ochrona narzędzi audytowych}\label{ochrona-narzedzi-audytowych}

Zgodnie z kontrolą \textbf{A.8.34 (Ochrona systemów informatycznych podczas testów audytowych)}:
\begin{itemize}
\tightlist
\item
  Narzędzia do skanowania podatności i audytu (np. Nessus, OpenVAS, nmap) zainstalowane wyłącznie na wyznaczonych stacjach roboczych administratora IT.
\item
  Dostęp do narzędzi audytowych chroniony hasłem i MFA.
\item
  Wyniki audytów i skanów klasyfikowane jako POUFNE - przechowywane w zaszyfrowanym repozytorium.
\item
  Audyty zewnętrzne (pen-test, przegląd ISMS) realizowane wyłącznie przez podmioty z podpisaną umową NDA i umową powierzenia przetwarzania danych (jeśli dotyczy).
\item
  Plan audytu zewnętrznego uzgadniany z wyprzedzeniem z właścicielem firmy.
\end{itemize}

\subsection{Raportowanie i przegląd kierowniczy}\label{raportowanie-i-przeglad-kierowniczy}

Zgodnie z wymaganiami sekcji 9.3 normy ISO/IEC 27001:2022 (Przegląd zarządzania):
\begin{itemize}
\tightlist
\item
  \textbf{Raport miesięczny} (administrator IT $\rightarrow$ właściciel firmy): podsumowanie alertów, wykrytych anomalii, statusu łatek bezpieczeństwa, wyników przeglądów logów.
\item
  \textbf{Raport kwartalny:} podsumowanie KPI bezpieczeństwa (MTTD, MTTR, wyniki testów phishingowych, status backupów), wyniki audytów uprawnień, wnioski z testów DRP.
\item
  \textbf{Raport roczny:} pełne podsumowanie stanu bezpieczeństwa, wyniki pen-testu, wyniki audytu RODO, ocena ryzyka, rekomendacje na kolejny rok.
\item
  Raporty przechowywane przez minimum 3 lata, klasyfikowane jako POUFNE.
\end{itemize}

\subsection{Rejestr zmian i przeglądów polityki}\label{rejestr-zmian-i-przegladow-polityki}

\begin{longtable}{@{}>{\raggedright\arraybackslash}p{\dimexpr\linewidth/1000*150-2\tabcolsep\relax}>{\raggedright\arraybackslash}p{\dimexpr\linewidth/1000*200-2\tabcolsep\relax}>{\raggedright\arraybackslash}p{\dimexpr\linewidth/1000*450-2\tabcolsep\relax}>{\raggedright\arraybackslash}p{\dimexpr\linewidth/1000*200-2\tabcolsep\relax}@{}}
\toprule\noalign{}
Wersja & Data & Opis zmiany & Autor \\
\midrule\noalign{}
\endhead
\bottomrule\noalign{}
\endlastfoot
1.0 & {} & Wersja początkowa (Dostosowanie do ISO/IEC 27001:2022) & {} \\
\end{longtable}

\begin{longtable}{@{}
  >{\raggedright\arraybackslash}p{\dimexpr\linewidth/1000*500-1\tabcolsep\relax}
  >{\raggedright\arraybackslash}p{\dimexpr\linewidth/1000*500-1\tabcolsep\relax}@{}}
\toprule\noalign{}
\endhead
\bottomrule\noalign{}
\endlastfoot
\begin{minipage}[t]{\linewidth}\raggedright
Zatwierdził:

Podpis: \underline{\hspace{3cm}}

Data: \underline{\hspace{4cm}}
\end{minipage} & \begin{minipage}[t]{\linewidth}\raggedright
Stanowisko / Tytuł:

Imię i nazwisko: \underline{\hspace{3cm}}

Funkcja: \underline{\hspace{3cm}}
\end{minipage} \\
\end{longtable}

\phantomsection\label{doc5}
PLAY-009

\section{PLAYBOOK SYTUACJI NIETYPOWYCH 
\newline (INCIDENT
RESPONSE) PLAY-009}\label{playbook-sytuacji-nietypowych-incident-response}

\begin{longtable}[]{@{}>{\raggedright\arraybackslash}p{\dimexpr\linewidth/1000*300-2\tabcolsep\relax}>{\raggedright\arraybackslash}p{\dimexpr\linewidth/1000*700-2\tabcolsep\relax}@{}}
\toprule\noalign{}
\endhead
\bottomrule\noalign{}
\endlastfoot
Wersja: & 1.0 \\
Dotyczy: & Wszyscy pracownicy \\
Odpowiedzialny: & Właściciel / Administrator IT \\
Aktualizacja: & Po każdym incydencie i co 12 miesięcy \\
\end{longtable}

\textbf{ALARM:} Ten dokument musi być dostępny OFFLINE (wydrukowany) w widocznym
miejscu. W przypadku incydentu sieć może być niedostępna.

\subsection{Kontakty alarmowe -- WYPEŁNIĆ I
WYDRUKOWAĆ}\label{kontakty-alarmowe-wypeux142niux107-i-wydrukowaux107}

\begin{longtable}[]{@{}>{\raggedright\arraybackslash}p{\dimexpr\linewidth/1000*270-2\tabcolsep\relax}>{\raggedright\arraybackslash}p{\dimexpr\linewidth/1000*250-2\tabcolsep\relax}>{\raggedright\arraybackslash}p{\dimexpr\linewidth/1000*250-2\tabcolsep\relax}>{\raggedright\arraybackslash}p{\dimexpr\linewidth/1000*230-2\tabcolsep\relax}@{}}
\toprule\noalign{}
Rola & Imię i Nazwisko & Telefon / E-mail & Zakres odpowiedzialności \\
\midrule\noalign{}
\endhead
\bottomrule\noalign{}
\endlastfoot
Właściciel firmy (CISO) & \underline{\hspace{3cm}} &
\underline{\hspace{3cm}} & Główna decyzja, autoryzacja działań \\
Administrator IT & \underline{\hspace{3cm}} &
\underline{\hspace{3cm}} & Techniczne działanie i izolacja \\
Osoba do kontaktu z Policją & \underline{\hspace{3cm}} &
\underline{\hspace{3cm}} & Zgłoszenia prawne, przestępstwa \\
Osoba do kontaktu z mediami & \underline{\hspace{3cm}} &
\underline{\hspace{3cm}} & JEDEN głos firmy, komunikacja zewn. \\
Osoba do kontaktu z UODO & \underline{\hspace{3cm}} &
\underline{\hspace{3cm}} & Naruszenia RODO (72h) \\
Prawnik firmy & \underline{\hspace{3cm}} &
\underline{\hspace{3cm}} & Porady prawne, kontakty z organami \\
Ubezpieczyciel (cyber) & \underline{\hspace{3cm}} &
\underline{\hspace{3cm}} & Zgłoszenie szkody \\
Zewnętrzna firma IT (backup) & \underline{\hspace{3cm}} &
\underline{\hspace{3cm}} & Wsparcie techniczne awaryjne \\
\end{longtable}

Ważne numery: Policja: 997 \textbar{} CERT Polska: cert.pl \textbar{}
UODO: uodo.gov.pl \textbar{} tel. UODO: 22 531 03 00

\subsection{Fazy reagowania na
incydent}\label{fazy-reagowania-na-incydent}

\subsubsection*{FAZA 1: Identyfikacja (0--30
min)}\label{faza-1-identyfikacja-030-min}

\begin{enumerate}
\tightlist
\item
  Pracownik wykrywa anomalię → \textbf{NATYCHMIASTOWO} informuje
  administratora IT i właściciela
\item
  Administrator IT ocenia powagę incydentu wg klasyfikacji poniżej
\item
  Rejestracja incydentu w Rejestrze Incydentów (godzina, opis, reporter)
\end{enumerate}

\subsubsection*{FAZA 2: Izolacja (30--60
min)}\label{faza-2-izolacja-3060-min}

\begin{enumerate}
\tightlist
\item
  Odizolowanie dotkniętych systemów od sieci (jeśli dotyczy) -- kabel
  LAN + WiFi
\item
  Zmiana haseł do zagrożonych kont
\item
  Zabezpieczenie dowodów (logi, kopie plików) -- \textbf{NIE USUWAĆ
  niczego}
\end{enumerate}

\subsubsection*{FAZA 3: Ograniczenie i
likwidacja}\label{faza-3-ograniczenie-i-likwidacja}

\begin{enumerate}
\tightlist
\item
  Blokada ruchu sieciowego do/ze źródeł ataku na firewall
\item
  Przywracanie z kopii zapasowej (wg DRP-011)
\item
  Identyfikacja wektora ataku i usunięcie złośliwego kodu
\end{enumerate}

\subsection{Matryca powiadamiania}\label{matryca-powiadamiania}

\begin{longtable}[]{@{}>{\raggedright\arraybackslash}p{\dimexpr\linewidth/1000*270-2\tabcolsep\relax}>{\raggedright\arraybackslash}p{\dimexpr\linewidth/1000*250-2\tabcolsep\relax}>{\raggedright\arraybackslash}p{\dimexpr\linewidth/1000*250-2\tabcolsep\relax}>{\raggedright\arraybackslash}p{\dimexpr\linewidth/1000*230-2\tabcolsep\relax}@{}}
\toprule\noalign{}
Kogo powiadomić & Kiedy & Kto powiadamia & Termin \\
\midrule\noalign{}
\endhead
\bottomrule\noalign{}
\endlastfoot
UODO (Urząd Ochrony Danych) & Naruszenie danych osobowych z ryzykiem &
Właściciel / IOD & \textbf{72 godziny} \\
Klienci poszkodowani & Wysokie ryzyko dla ich danych & Właściciel & Bez
zwłoki \\
Policja / CERT Polska & Przestępstwo, atak hakerski, ransomware &
Wyznaczona osoba & Niezwłocznie \\
Media / PR & Jeśli incydent jest/może być publiczny & Właściciel (JEDEN
głos) & Po konsultacji z prawnikiem \\
Ubezpieczyciel & Szkoda objęta polisą cyber & Właściciel & 24--48
godzin \\
\end{longtable}

\subsection{Scenariusze -- reakcje}\label{scenariusze-reakcje}

\subsubsection*{Scenariusz A: Phishing / kompromitacja
konta}\label{scenariusz-a-phishing-kompromitacja-konta}

\begin{itemize}
\tightlist
\item
  Natychmiast zmień hasło i zresetuj MFA dla konta
\item
  Przegląd aktywności konta (logi) za ostatnie 30 dni
\item
  Sprawdź czy dane były eksportowane lub przesyłane
\item
  Poinformuj pozostałych pracowników (alert phishingowy)
\end{itemize}

\subsubsection*{Scenariusz B:
Ransomware}\label{scenariusz-b-ransomware}

\begin{itemize}
\tightlist
\item
  \textbf{NATYCHMIASTOWO} odłącz zainfekowane urządzenie od sieci (kabel
  + WiFi)
\item
  \textbf{NIE PŁACIĆ OKUPU} bez konsultacji z prawnikiem i
  ubezpieczycielem
\item
  Uruchom procedurę DRP-011
\item
  Zgłoś do Policji i CERT Polska (cert.pl)
\item
  Zidentyfikuj wektor ataku i załataj podatność
\end{itemize}

\subsubsection*{Scenariusz C: Wyciek danych
klientów}\label{scenariusz-c-wyciek-danych-klientuxf3w}

\begin{itemize}
\tightlist
\item
  Zidentyfikuj jakie dane wyciekły i ile osób dotyczy
\item
  Ocena ryzyka wg RODO -- czy zgłosić do UODO (72h)?
\item
  Powiadom poszkodowanych klientów
\item
  Zabezpiecz dowody (nie usuwaj logów)
\end{itemize}

\subsubsection*{Scenariusz D: Fraud finansowy / zmiana rachunku}\label{scenariusz-D-kompromitacja-systemu}
\begin{itemize}
\tightlist
\item
\textbf{NATYCHMIASTOWO} wstrzymaj przelew
\item
Skontaktować się z bankiem
\item
Skontaktowanie się z klientem
\item
Zabezpiecz korespondencję e-mail
\item
Zgłoś incydent do Policji
\item 
Zarejestruj incydent w Rejestrze Incydent
\end{itemize}

\begin{longtable}[]{@{}>{\raggedright\arraybackslash}p{\dimexpr\linewidth/1000*220-2\tabcolsep\relax}>{\raggedright\arraybackslash}p{\dimexpr\linewidth/1000*200-2\tabcolsep\relax}>{\raggedright\arraybackslash}p{\dimexpr\linewidth/1000*200-2\tabcolsep\relax}>{\raggedright\arraybackslash}p{\dimexpr\linewidth/1000*200-2\tabcolsep\relax}>{\raggedright\arraybackslash}p{\dimexpr\linewidth/1000*180-2\tabcolsep\relax}@{}}
\toprule\noalign{}
Data/Czas & Opis incydentu & Działania podjęte & Status & Zamknięty \\
\midrule\noalign{}
\endhead
\bottomrule\noalign{}
\endlastfoot
\underline{\hspace{2cm}} &
\underline{\hspace{2cm}} &
\underline{\hspace{2cm}} & {[} {]} Otwarty &
\underline{\hspace{2cm}} \\
\underline{\hspace{2cm}} &
\underline{\hspace{2cm}} &
\underline{\hspace{2cm}} & {[} {]} Otwarty &
\underline{\hspace{2cm}} \\
\end{longtable}

\begin{longtable}[]{@{}
  >{\raggedright\arraybackslash}p{\dimexpr\linewidth/1000*500-1\tabcolsep\relax}
  >{\raggedright\arraybackslash}p{\dimexpr\linewidth/1000*500-1\tabcolsep\relax}@{}}
\toprule\noalign{}
\endhead
\bottomrule\noalign{}
\endlastfoot
\begin{minipage}[t]{\linewidth}\raggedright
Zatwierdził:

Podpis: \underline{\hspace{3cm}}

Data: \underline{\hspace{4cm}}
\end{minipage} & \begin{minipage}[t]{\linewidth}\raggedright
Stanowisko / Tytuł:

Imię i nazwisko: \underline{\hspace{3cm}}

Funkcja: \underline{\hspace{3cm}}
\end{minipage} \\
\end{longtable}

\phantomsection\label{doc8}
DRP-010

\section{PLAN ODTWARZANIA PO KATASTROFIE 
\newline (DISASTER RECOVERY PLAN) DRP-010}\label{plan-odtwarzania-po-katastrofie-disaster-recovery-plan}

\begin{longtable}[]{@{}>{\raggedright\arraybackslash}p{\dimexpr\linewidth/1000*350-2\tabcolsep\relax}>{\raggedright\arraybackslash}p{\dimexpr\linewidth/1000*650-2\tabcolsep\relax}@{}}
\toprule\noalign{}
\endhead
\bottomrule\noalign{}
\endlastfoot
Wersja: & 1.0 \\
Odpowiedzialny: & Właściciel / Administrator IT \\
Globalne RTO (dla procesów biznesowych): & 19-20 oraz 24-25 dzień miesiąca: 6 godzin \newline Pozostałe dni miesiąca: 24 godziny \\
RPO (dla danych księgowych): & 4 godziny \\
Testowanie planu: & Co 12 miesięcy (test pełny) + co 3 miesiące (test częściowy) \\
\end{longtable}

\subsection{Sztab Kryzysowy i deklaracja katastrofy}\label{sztab-kryzysowy-i-deklaracja-katastrofy}

Procedurę DRP (ogłoszenie stanu katastrofy) może uruchomić wyłącznie wyznaczona osoba ze Sztabu Kryzysowego, po wcześniejszej ocenie powagi incydentu. Samodzielne podjęcie działań odtworzeniowych przez pracowników niższego szczebla jest zabronione.

\begin{longtable}[]{@{}>{\raggedright\arraybackslash}p{\dimexpr\linewidth/1000*200-2\tabcolsep\relax}>{\raggedright\arraybackslash}p{\dimexpr\linewidth/1000*300-2\tabcolsep\relax}>{\raggedright\arraybackslash}p{\dimexpr\linewidth/1000*500-2\tabcolsep\relax}@{}}
\toprule\noalign{}
Rola w kryzysie & Stanowisko & Obowiązki / Uprawnienia \\
\midrule\noalign{}
\endhead
\bottomrule\noalign{}
\endlastfoot
Koordynator DRP & Właściciel / Zarząd & Podejmuje formalną decyzję o ogłoszeniu katastrofy. Zarządza komunikacją biznesową i prawną. \\
Główny Wykonawca & Główny Administrator IT & Koordynuje techniczne działania odtworzeniowe. Nadzoruje pracę zespołów obszarowych, weryfikuje poprawność przywracania systemów i raportuje do Zarządu. \\
Specjaliści ds. Odtwarzania & Zespoły IT (Administratorzy Obszarowi) & Realizują techniczne kroki DRP w przypisanych im obszarach (Infrastruktura sieciowa, Samba, Hosting, AD, VPN) wg ustalonej kolejności. Raportują postępy Głównemu Administratorowi. \\
Lider Biznesowy & Główny Księgowy & Weryfikuje integralność odzyskanych danych (np. po przywróceniu systemu księgowego). Informuje zespół księgowy o gotowości środowiska do dalszej pracy. \\
\end{longtable}

\subsection{Zapasowy kanał komunikacji}\label{zapasowy-kanal-komunikacji}

W przypadku krytycznej awarii infrastruktury IT (np. niedostępność kontrolera domeny Active Directory, poczty e-mail czy infrastruktury sieciowej), standardowe firmowe kanały komunikacji mogą być całkowicie niedostępne. W celu zapewnienia sprawnego zarządzania incydentem, ustala się następujące zasady komunikacji awaryjnej:

\begin{itemize}
\tightlist
\item
  \textbf{Powiadamianie pracowników (kanał ogólny):} Głównym zapasowym kanałem informowania pracowników firmy księgowej o awarii, przewidywanym czasie przestoju (RTO) oraz o wznowieniu pracy są \textbf{zbiorcze komunikaty SMS}. Wysyłką SMS-ów zarządza Koordynator DRP (lub wyznaczony Lider Biznesowy).
\item
  \textbf{Komunikacja operacyjna (Sztab Kryzysowy i zespoły IT):} Do bieżącego koordynowania prac naprawczych pomiędzy Administratorami (odpowiedzialnymi za sieć, Sambę, AD, hosting) wykorzystywane są bezpośrednie połączenia telefoniczne.
\item
  \textbf{Awaryjna baza kontaktów:} Aby komunikacja SMS była możliwa w sytuacji braku dostępu do serwerów i bazy danych, \textbf{wydrukowana kopia} aktualnej listy numerów telefonów do wszystkich pracowników, członków Sztabu Kryzysowego oraz kluczowych dostawców zewnętrznych (np. dostawca Hostingu) musi być przechowywana w bezpiecznym, fizycznym miejscu (np. sejf w biurze Właściciela). Lista ta jest aktualizowana co 6 miesięcy.
\end{itemize}

\subsection{Klasyfikacja systemów
krytycznych}\label{klasyfikacja-systemuxf3w-krytycznych}

\begin{longtable}[]{@{}>{\raggedright\arraybackslash}p{\dimexpr\linewidth/1000*270-2\tabcolsep\relax}>{\raggedright\arraybackslash}p{\dimexpr\linewidth/1000*250-2\tabcolsep\relax}>{\raggedright\arraybackslash}p{\dimexpr\linewidth/1000*250-2\tabcolsep\relax}>{\raggedright\arraybackslash}p{\dimexpr\linewidth/1000*230-2\tabcolsep\relax}@{}}
\toprule\noalign{}
System & Priorytet & RTO (maks.) & RPO (maks.) \\
\midrule\noalign{}
\endhead
\bottomrule\noalign{}
\endlastfoot
System księgowy (główna aplikacja) & KRYTYCZNY (P1) & 2 godziny & 4
godziny \\
Baza danych klientów & KRYTYCZNY (P1) & 2 godziny & 4 godziny \\
Windows Server / Active Directory & KRYTYCZNY (P1) & 2 godziny & 4 godziny \\
Poczta e-mail & WYSOKI (P2) & 4 godziny & 24 godziny \\
Przestrzeń dyskowa (Samba) & WYSOKI (P2) & 4 godziny & 24
godziny \\
VPN / Zdalny dostęp & ŚREDNI (P3) & 8 godzin & 48 godzin \\
Hosting (Strona WWW) & NISKI (P4) & 48 godzin & 7 dni \\
\end{longtable}

\subsection{Kolejność przywracania
systemów}\label{kolejnoux15bux107-przywracania-systemuxf3w}

\begin{enumerate}
\tightlist
\item
  \textbf{BEZPIECZEŃSTWO} -- czy środowisko jest bezpieczne? (brak
  aktywnego ataku)
\item
  Infrastruktura sieciowa (router, switch, firewall) -- konfiguracja z
  backupu
\item
  Windows Server / Active Directory -- zarządzanie tożsamością
\item
  Serwer baz danych -- przywracanie z ostatniego czystego backupu
\item
  System księgowy -- weryfikacja integralności danych
\item
  Przestrzeń dyskowa (Samba) -- przywrócenie plików i struktury katalogów z backupu
\item
  Weryfikacja poprawnego logowania i dostępu do udziałów (Samby)
\item
  Poczta e-mail
\item
  VPN i dostęp zdalny
\item
  Weryfikacja kompletności -- porównanie z RPO
\item
  Informacja dla pracowników o wznowieniu pracy
\item
  Dokumentacja incydentu, wnioski (lessons learned), aktualizacja DRP
\end{enumerate}

\subsection{Procedura powrotu do infrastruktury głównej (Failback)}\label{procedura-powrotu-failback}

Po zażegnaniu kryzysu, naprawieniu głównej infrastruktury lub wdrożeniu nowego sprzętu, następuje zaplanowany powrót do standardowego trybu pracy (tzw. Failback). Proces ten jest krytyczny, aby nie utracić danych wprowadzonych przez księgowe w czasie działania na środowisku awaryjnym.

\newpage
\textbf{Główne kroki procedury Failback:}
\begin{enumerate}
\tightlist
\item
  \textbf{Weryfikacja docelowego środowiska:} Główny Administrator IT potwierdza pełną sprawność, stabilność i bezpieczeństwo (brak podatności/malware) naprawionej infrastruktury (Active Directory, sieć, hosting).
\item
  \textbf{Zaplanowanie okna serwisowego:} Ustalenie z Liderem Biznesowym (Głównym Księgowym) terminu przełączenia. Przepięcie musi nastąpić poza godzinami pracy biura, aby uniknąć rozbieżności w bazach danych.
\item
  \textbf{Synchronizacja danych (Reverse Synchronization):} Wykonanie kopii zapasowej aktualnych danych ze środowiska awaryjnego (baza danych klientów, pliki na serwerze Samba) i odtworzenie ich na infrastrukturze głównej.
\item
  \textbf{Przepięcie usług sieciowych:} Przekierowanie ruchu z powrotem na główne serwery (np. rekonfiguracja routingu dla VPN, aktualizacja rekordów DNS dla Hostingu).
\item
  \textbf{Testy funkcjonalne i wznowienie pracy:} Potwierdzenie dostępu do Active Directory i zasobów, a następnie poinformowanie pracowników o zakończeniu okna serwisowego.
\item
  \textbf{Zakończenie DRP:} Odtworzenie standardowych harmonogramów kopii zapasowych dla głównej infrastruktury i oficjalne zamknięcie incydentu z wpisem do dokumentacji (Lessons Learned).
\end{enumerate}



\subsection{Lokalizacje kopii zapasowych}\label{lokalizacje-kopii-zapasowych}

\begin{longtable}[]{@{}>{\raggedright\arraybackslash}p{\dimexpr\linewidth/1000*270-2\tabcolsep\relax}>{\raggedright\arraybackslash}p{\dimexpr\linewidth/1000*250-2\tabcolsep\relax}>{\raggedright\arraybackslash}p{\dimexpr\linewidth/1000*250-2\tabcolsep\relax}>{\raggedright\arraybackslash}p{\dimexpr\linewidth/1000*230-2\tabcolsep\relax}@{}}
\toprule\noalign{}
Rodzaj kopii & Lokalizacja & Dostęp (kto / jak) & Przetestowany \\
\midrule\noalign{}
\endhead
\bottomrule\noalign{}
\endlastfoot
Backup lokalny (NAS) & {[}LOKALIZACJA{]} & {[}DANE DOSTĘPU{]} & {[} {]}
\underline{\hspace{2cm}} \\
Backup offsite (fizyczny) & {[}LOKALIZACJA{]} & {[}OSOBA
ODPOWIEDZIALNA{]} & {[} {]} \underline{\hspace{2cm}} \\
Backup chmurowy & {[}PLATFORMA/URL{]} & {[}DANE DOSTĘPU{]} & {[} {]}
\underline{\hspace{2cm}} \\
\end{longtable}

\subsection{Procedura testowania DRP}\label{procedura-testowania-drp}

\begin{itemize}
\tightlist
\item
  \textbf{Test częściowy (co kwartał):} przywracanie wybranego systemu
  lub pliku z backupów -- dokumentacja czasu i wyników
\item
  \textbf{Test pełny (co 12 miesięcy):} symulacja pełnej utraty
  środowiska (możliwe środowisko testowe)
\item
  Dokumentacja wyników testu: osiągnięty RTO, RPO, uchybienia
\item
  Aktualizacja DRP na podstawie wyników testów
\end{itemize}

\begin{longtable}[]{@{}
  >{\raggedright\arraybackslash}p{\dimexpr\linewidth/1000*500-1\tabcolsep\relax}
  >{\raggedright\arraybackslash}p{\dimexpr\linewidth/1000*500-1\tabcolsep\relax}@{}}
\toprule\noalign{}
\endhead
\bottomrule\noalign{}
\endlastfoot
\begin{minipage}[t]{\linewidth}\raggedright
Zatwierdził:

Podpis: \underline{\hspace{3cm}}

Data: \underline{\hspace{4cm}}
\end{minipage} & \begin{minipage}[t]{\linewidth}\raggedright
Stanowisko / Tytuł:

Imię i nazwisko: \underline{\hspace{3cm}}

Funkcja: \underline{\hspace{3cm}}
\end{minipage} \\
\end{longtable}

\phantomsection\label{doc12}
PLAN-011

\section{PLAN CIĄGŁEGO DOSKONALENIA
CYBERBEZPIECZEŃSTWA PLAN-011}\label{plan-ciux105gux142ego-doskonalenia-cyberbezpieczeux144stwa}

\begin{longtable}[]{@{}>{\raggedright\arraybackslash}p{\dimexpr\linewidth/1000*300-2\tabcolsep\relax}>{\raggedright\arraybackslash}p{\dimexpr\linewidth/1000*700-2\tabcolsep\relax}@{}}
\toprule\noalign{}
\endhead
\bottomrule\noalign{}
\endlastfoot
Wersja: & 1.0 \\
Odpowiedzialny: & Właściciel / Administrator IT \\
Cykl przeglądu: & Kwartalny przegląd KPI + roczne podsumowanie i
planowanie \\
\end{longtable}

Cyberbezpieczeństwo nie jest jednorazowym projektem, lecz procesem.
Firma zobowiązuje się do regularnego przeglądu, doskonalenia i
rozwijania swoich praktyk bezpieczeństwa.

\subsection{Cykl PDCA dla cyberbezpieczeństwa}\label{cykl-pdca-dla-cyberbezpieczeux144stwa}

\begin{longtable}[]{@{}>{\raggedright\arraybackslash}p{\dimexpr\linewidth/1000*270-2\tabcolsep\relax}>{\raggedright\arraybackslash}p{\dimexpr\linewidth/1000*430-2\tabcolsep\relax}>{\raggedright\arraybackslash}p{\dimexpr\linewidth/1000*300-2\tabcolsep\relax}@{}}
\toprule\noalign{}
Faza & Działanie & Przykłady aktywności \\
\midrule\noalign{}
\endhead
\bottomrule\noalign{}
\endlastfoot
{[P]} PLAN (Planowanie) & Ocena ryzyka, ustalenie celów & Roczna ocena ryzyka, ustalenie KPI bezpieczeństwa \\
{[D]} DO (Realizacja) & Wdrożenie środków & Nowe narzędzia, szkolenia, zaktualizowane procedury \\
{[C]} CHECK (Sprawdzenie) & Monitorowanie i audyt & Testy penetracyjne, przeglądy logów, KPI \\
{[A]} ACT (Działanie) & Korekty i ulepszenia & Aktualizacja polityk po incydencie / audycie \\
\end{longtable}

\newpage
\subsection{Roczna ocena ryzyka
(szablon)}\label{roczna-ocena-ryzyka-szablon}
\begin{longtable}[]{@{}>{\raggedright\arraybackslash}p{\dimexpr\linewidth/1000*270-2\tabcolsep\relax}>{\raggedright\arraybackslash}p{\dimexpr\linewidth/1000*200-2\tabcolsep\relax}>{\raggedright\arraybackslash}p{\dimexpr\linewidth/1000*150-2\tabcolsep\relax}>{\raggedright\arraybackslash}p{\dimexpr\linewidth/1000*380-2\tabcolsep\relax}@{}}
\toprule\noalign{}
Ryzyko & Prawdopodo\-bieństwo \newline (1--5) & Wpływ \newline (1--5) & Planowane działanie (Korekcyjne/Zapobiegawcze) \\
%Ryzyko & Prawdopodobieństwo (1--5) & Wpływ (1--5) & Planowane działanie (Korekcyjne/Zapobiegawcze) \\
\midrule\noalign{}
\endhead
\bottomrule\noalign{}
\endlastfoot
Phishing / kompromitacja konta Active Directory & 4 & 5 & Wdrożenie MFA (Multi-Factor Authentication) dla wszystkich kont, kwartalne szkolenia z cyberbezpieczeństwa. \\
Ransomware (szyfrowanie udziałów Samba i baz danych) & 3 & 5 & Segmentacja sieci, zaostrzenie uprawnień na Sambie (zasada najmniejszych uprawnień), weryfikacja backupów offline. \\
Wyciek danych klientów (podatność w Hostingu/WWW) & 2 & 5 & Kwartalne skanowanie podatności aplikacji webowej, wdrożenie zapory WAF (Web Application Firewall). \\
Nieautoryzowany dostęp z zewnątrz (VPN brute-force) & 3 & 4 & Wymuszenie silnych haseł w AD, wdrożenie blokady konta po 5 błędnych próbach, geoblokowanie ruchu na VPN. \\
Utrata backupów / Awaria sprzętowa NAS & 2 & 4 & Wdrożenie reguły backupu 3-2-1 (w tym kopia w chmurze lub offline wg. procedur DRP). \\
Awaria głównej infrastruktury sieciowej (Router/Switch) & 2 & 3 & Przechowywanie zapasowych konfiguracji (backup config) na bezpiecznym serwerze, umowa SLA z dostawcą sprzętu. \\
\end{longtable}

\subsection{KPI -- wskaźniki poziomu
bezpieczeństwa}\label{kpi-wskaux17aniki-poziomu-bezpieczeux144stwa}

\begin{longtable}[]{@{}>{\raggedright\arraybackslash}p{\dimexpr\linewidth/1000*320-2\tabcolsep\relax}>{\raggedright\arraybackslash}p{\dimexpr\linewidth/1000*200-2\tabcolsep\relax}>{\raggedright\arraybackslash}p{\dimexpr\linewidth/1000*250-2\tabcolsep\relax}>{\raggedright\arraybackslash}p{\dimexpr\linewidth/1000*230-2\tabcolsep\relax}@{}}
\toprule\noalign{}
Wskaźnik & Cel & Aktualny wynik & Trend \\
\midrule\noalign{}
\endhead
\bottomrule\noalign{}
\endlastfoot
Czas wdrożenia krytycznych łat (Windows Server / VPN) & \textless{} 72 godziny od publikacji & \underline{\hspace{2cm}} & {[} {]} ↑ {[} {]} ↓ \\
Testy przywracania systemów DRP (AD, Samba, Baza) & 100\% zakończonych sukcesem & \underline{\hspace{2cm}} & {[} {]} ↑ {[} {]} ↓ \\
Czas średniego reagowania na incydent (MTTR) & \textless{} 4 godziny & \underline{\hspace{2cm}} & {[} {]} ↑ {[} {]} ↓ \\
\% pracowników z aktualnym szkoleniem Security Awareness & 100\% & \underline{\hspace{2cm}} & {[} {]} ↑ {[} {]} ↓ \\
Skuteczne logowania VPN bez użycia MFA & 0 & \underline{\hspace{2cm}} & {[} {]} ↑ {[} {]} ↓ \\
Współczynnik kliknięć w symulowanych testach Phishingowych & \textless{} 5\% & \underline{\hspace{2cm}} & {[} {]} ↑ {[} {]} ↓ \\
\end{longtable}

\subsection{Harmonogram działań
doskonalenia}\label{harmonogram-dziaux142aux144-doskonalenia}

\begin{longtable}[]{@{}>{\raggedright\arraybackslash}p{\dimexpr\linewidth/1000*320-2\tabcolsep\relax}>{\raggedright\arraybackslash}p{\dimexpr\linewidth/1000*200-2\tabcolsep\relax}>{\raggedright\arraybackslash}p{\dimexpr\linewidth/1000*250-2\tabcolsep\relax}>{\raggedright\arraybackslash}p{\dimexpr\linewidth/1000*230-2\tabcolsep\relax}@{}}
\toprule\noalign{}
Działanie & Termin & Odpowiedzialny & Status \\
\midrule\noalign{}
\endhead
\bottomrule\noalign{}
\endlastfoot
Roczna ocena ryzyka dla systemów IT & Styczeń & Właściciel / Główny Admin IT & {[} {]} \underline{\hspace{2cm}} \\
Automatyczne skanowanie podatności (VPN, Hosting) & Co kwartał & Zespół Sieci / Hostingu & {[} {]} \underline{\hspace{2cm}} \\
Wewnętrzny audyt bezpieczeństwa (Zgodność z ISO/RODO) & Q3 & Wyznaczony Audytor & {[} {]} \underline{\hspace{2cm}} \\
Test DRP (pełny dla AD, Samby i Baz Danych) & Q4 & Główny Admin IT & {[} {]} \underline{\hspace{2cm}} \\
Szkolenia odświeżające (Anty-phishing) & Co 12 miesięcy & Dział IT & {[} {]} \underline{\hspace{2cm}} \\
Przegląd uprawnień użytkowników w Active Directory i Sambie & Co kwartał & Zespół AD / Samba & {[} {]} \underline{\hspace{2cm}} \\
Przegląd Zarządzania bezpieczeństwem informacji (Management Review) & Grudzień (Q4) & Właściciel / Zarząd & {[} {]}\underline{\hspace{2cm}} \\
\end{longtable}

\begin{longtable}[]{@{}
  >{\raggedright\arraybackslash}p{\dimexpr\linewidth/1000*500-1\tabcolsep\relax}
  >{\raggedright\arraybackslash}p{\dimexpr\linewidth/1000*500-1\tabcolsep\relax}@{}}
\toprule\noalign{}
\endhead
\bottomrule\noalign{}
\endlastfoot
\begin{minipage}[t]{\linewidth}\raggedright
Zatwierdził:

Podpis: \underline{\hspace{3cm}}

Data: \underline{\hspace{4cm}}
\end{minipage} & \begin{minipage}[t]{\linewidth}\raggedright
Stanowisko / Tytuł:

Imię i nazwisko: \underline{\hspace{3cm}}

Funkcja: \underline{\hspace{3cm}}
\end{minipage} \\
\end{longtable}

\phantomsection\label{doc12}
PROC-012

\section{PROCEDURA SZKOLEŃ I WERYFIKACJI WIEDZY Z
CYBERBEZPIECZEŃSTWA PROC-012}\label{procedura-szkoleux144-i-weryfikacji-wiedzy-z-cyberbezpieczeux144stwa}

\begin{longtable}[]{@{}>{\raggedright\arraybackslash}p{\dimexpr\linewidth/1000*300-2\tabcolsep\relax}>{\raggedright\arraybackslash}p{\dimexpr\linewidth/1000*700-2\tabcolsep\relax}@{}}
\toprule\noalign{}
\endhead
\bottomrule\noalign{}
\endlastfoot
Wersja: & 1.0 \\
Dotyczy: & Wszyscy pracownicy i współpracownicy \\
Odpowiedzialny: & Administrator IT + Właściciel \\
Częstotliwość: & Szkolenie bazowe: onboarding \textbar{} Odświeżające:
co 12 miesięcy \\
\end{longtable}

\subsection{Program szkoleń}\label{program-szkoleux144}

\subsubsection{Szkolenie bazowe (nowi
pracownicy)}\label{szkolenie-bazowe-nowi-pracownicy}

\begin{itemize}
\tightlist
\item
  Czas trwania: min. 2 godziny
\item
  Format: prezentacja + omówienie + quiz
\item
  Zawartość: polityki IT, hasła i MFA, phishing i inżynieria społeczna,
  obsługa incydentów, RODO, bezpieczna praca zdalna
  Dodatkowy moduł dla pracowników księgowości zawierający:
  obsługę faktur, przelewów, deklaracji podatkowych oraz komunikacji z klientami.
\item
  Wymagany wynik quizu: min. \textbf{80\%}
\item
  Powtarzanie przy wyniku poniżej 80\% (max 3 podejścia, potem szkolenie
  indywidualne)
\end{itemize}

\subsubsection{Szkolenie odświeżające (co 12
miesięcy)}\label{szkolenie-odux15bwieux17cajux105ce-co-12-miesiux119cy}

\begin{itemize}
\tightlist
\item
  Czas trwania: min. 1 godzina
\item
  Format: e-learning lub warsztat stacjonarny
\item
  Zawartość: nowe zagrożenia roku, przegląd incydentów z minionego roku,
  nowe procedury.
\item
  Wymagany wynik testu: min. \textbf{80\%}
\end{itemize}

\subsection{Testy phishingowe}\label{testy-phishingowe}

Firma przeprowadza symulowane ataki phishingowe w celu weryfikacji
świadomości pracowników:

\begin{longtable}[]{@{}>{\raggedright\arraybackslash}p{\dimexpr\linewidth/1000*270-2\tabcolsep\relax}>{\raggedright\arraybackslash}p{\dimexpr\linewidth/1000*430-2\tabcolsep\relax}>{\raggedright\arraybackslash}p{\dimexpr\linewidth/1000*300-2\tabcolsep\relax}@{}}
\toprule\noalign{}
Parametr & Opis & Wartość docelowa \\
\midrule\noalign{}
\endhead
\bottomrule\noalign{}
\endlastfoot
Częstotliwość & Niezapowiedziane testy & Co kwartał (minimum) \\
Narzędzie & Platforma do symulacji phishingu & GoPhish \\
Próg alertu & Odsetek kliknięć w link & Cel: \textless10\%
pracowników \\
Reakcja na kliknięcie & Natychmiastowe szkolenie remediacyjne &
Obowiązkowe w ciągu 7 dni \\
Raportowanie & Anonimowe wyniki zbiorcze & Co kwartał -- Administrator IT \\
\end{longtable}

\subsection{Przegląd Playbooków i
ćwiczenia}\label{przeglux105d-playbookuxf3w-i-ux107wiczenia}

\begin{itemize}
\tightlist
\item
  Pracownicy zapoznają się z Playbookiem Incident Response (PLAY-009) co
  12 miesięcy
\item
  Coroczne ćwiczenie tablicowe (tabletop exercise) -- symulacja
  scenariusza incydentu
\item
  Udokumentowane wyniki ćwiczenia i wnioski (lessons learned)
\end{itemize}

\subsection{Rejestr szkoleń (szablon) }\label{rejestr-szkoleux144}

\begin{longtable}[]{@{}>{\raggedright\arraybackslash}p{\dimexpr\linewidth/1000*220-2\tabcolsep\relax}>{\raggedright\arraybackslash}p{\dimexpr\linewidth/1000*200-2\tabcolsep\relax}>{\raggedright\arraybackslash}p{\dimexpr\linewidth/1000*200-2\tabcolsep\relax}>{\raggedright\arraybackslash}p{\dimexpr\linewidth/1000*200-2\tabcolsep\relax}>{\raggedright\arraybackslash}p{\dimexpr\linewidth/1000*180-2\tabcolsep\relax}@{}}
\toprule\noalign{}
Pracownik & Szkolenie & Data & Wynik testu & Certyfikat \\
\midrule\noalign{}
\endhead
\bottomrule\noalign{}
\endlastfoot
\underline{\hspace{3cm}} & Bazowe & \underline{\hspace{2cm}} &
\underline{\hspace{2cm}}\% & {[} {]} TAK \\
\underline{\hspace{3cm}} & Odświeżające & \underline{\hspace{2cm}} &
\underline{\hspace{2cm}}\% & {[} {]} TAK \\
\underline{\hspace{3cm}} & Phishing test & \underline{\hspace{2cm}}
& Kliknął: {[} {]} TAK / NIE & N/A \\
\underline{\hspace{3cm}} & Bazowe & \underline{\hspace{2cm}} &
\underline{\hspace{2cm}}\% & {[} {]} TAK \\
\end{longtable}

\begin{longtable}[]{@{}
  >{\raggedright\arraybackslash}p{\dimexpr\linewidth/1000*500-1\tabcolsep\relax}
  >{\raggedright\arraybackslash}p{\dimexpr\linewidth/1000*500-1\tabcolsep\relax}@{}}
\toprule\noalign{}
\endhead
\bottomrule\noalign{}
\endlastfoot
\begin{minipage}[t]{\linewidth}\raggedright
Zatwierdził:

Podpis: \underline{\hspace{3cm}}

Data: \underline{\hspace{4cm}}
\end{minipage} & \begin{minipage}[t]{\linewidth}\raggedright
Stanowisko / Tytuł:

Imię i nazwisko: \underline{\hspace{3cm}}

Funkcja: \underline{\hspace{3cm}}
\end{minipage} \\
\end{longtable}

{[}NAZWA FIRMY KSIĘGOWEJ{]} • Pakiet Dokumentacji Cyberbezpieczeństwa
v1.0 • POUFNE

\end{document}